% =====================================================================
%  SYMPHONY V2 - Manuale per il firmwarista (Engineering Model)
%  Riferimento: schematico rev. 1.45 del log hardware, 04/11/2026
% =====================================================================
\documentclass[11pt,a4paper,twoside,openany]{report}

% ---------------------------------------------------------------- base
\usepackage[T1]{fontenc}
\usepackage[utf8]{inputenc}
\usepackage[italian]{babel}
\usepackage{libertinus}
\usepackage[scaled=0.84]{beramono}
\usepackage{microtype}
\usepackage[a4paper,inner=26mm,outer=22mm,top=26mm,bottom=26mm,headheight=15pt]{geometry}
\usepackage[table,dvipsnames]{xcolor}
\usepackage{graphicx}
\usepackage{amsmath,amssymb}
\usepackage[autostyle]{csquotes}
\usepackage{enumitem}
\usepackage{pdflscape}
\usepackage{array,booktabs,longtable,xltabular,multirow,makecell}
\usepackage{siunitx}
\usepackage{fancyhdr}
\usepackage{titlesec}
\usepackage{tocloft}
\usepackage[most]{tcolorbox}
\usepackage{tikz}
\usetikzlibrary{arrows.meta,positioning,calc,fit,backgrounds,shapes.geometric,decorations.pathreplacing}
\usepackage[european,cuteinductors,straightvoltages]{circuitikz}
\usepackage{tikz-timing}
\usepackage{pgfplots}
\pgfplotsset{compat=1.18}
\usepackage{caption}
\usepackage{hyperref}
\usepackage{bookmark}

% ---------------------------------------------------------------- colori
\definecolor{sbblue}{HTML}{1F4E79}
\definecolor{sbblue2}{HTML}{2E75B6}
\definecolor{sblight}{HTML}{EAF1F8}
\definecolor{sbrow}{HTML}{F3F6FA}
\definecolor{sbred}{HTML}{B03A2E}
\definecolor{sbredbg}{HTML}{FBEDEB}
\definecolor{sbgreen}{HTML}{2E7D4F}
\definecolor{sbgreenbg}{HTML}{EAF5EE}
\definecolor{sbgray}{HTML}{5F6B76}
\definecolor{sbgraybg}{HTML}{F2F4F6}
\definecolor{sborange}{HTML}{C0730A}
\definecolor{sborangebg}{HTML}{FDF3E4}

\hypersetup{colorlinks=true,linkcolor=sbblue,urlcolor=sbblue2,citecolor=sbblue,
  pdftitle={SYMPHONY V2 - Manuale per il firmwarista},
  pdfauthor={Team SYMPHONY},pdfsubject={Engineering Model, schematico rev. 1.45}}

% ---------------------------------------------------------------- siunitx
\sisetup{output-decimal-marker={,},group-separator={\,},per-mode=symbol,
  range-phrase={\,\textendash\,},range-units=single,range-open-phrase={},list-final-separator={ e },
  detect-all}
\DeclareSIUnit\dBm{dBm}
\DeclareSIUnit\dBc{dBc}
\DeclareSIUnit\ppm{ppm}

% ---------------------------------------------------------------- titoli
\titleformat{\chapter}[display]
  {\normalfont\sffamily\color{sbblue}}
  {\begin{tikzpicture}[baseline]\node[fill=sbblue,text=white,minimum size=13mm,
     font=\sffamily\bfseries\Huge,inner sep=0]{\thechapter};\end{tikzpicture}}
  {12pt}{\Huge\bfseries}[\vspace{2pt}{\color{sbblue2}\titlerule[1.2pt]}]
\titlespacing*{\chapter}{0pt}{-10pt}{24pt}
\titleformat{\section}{\normalfont\sffamily\Large\bfseries\color{sbblue}}{\thesection}{0.8em}{}
\titleformat{\subsection}{\normalfont\sffamily\large\bfseries\color{sbblue2}}{\thesubsection}{0.8em}{}
\titleformat{\subsubsection}{\normalfont\sffamily\normalsize\bfseries\color{sbgray}}{\thesubsubsection}{0.8em}{}
\setcounter{secnumdepth}{3}
\setcounter{tocdepth}{1}
\renewcommand{\cftchapfont}{\sffamily\bfseries}
\renewcommand{\cftchappagefont}{\sffamily\bfseries}

% ---------------------------------------------------------------- testatine
\pagestyle{fancy}
\fancyhf{}
\renewcommand{\headrulewidth}{0.4pt}
\renewcommand{\headrule}{\hbox to\headwidth{\color{sbblue2}\leaders\hrule height \headrulewidth\hfill}}
\fancyhead[LE]{\sffamily\small\color{sbgray}SYMPHONY V2 \textbullet{} Manuale per il firmwarista}
\fancyhead[RO]{\sffamily\small\color{sbgray}\nouppercase{\leftmark}}
\fancyfoot[LE,RO]{\sffamily\small\bfseries\color{sbblue}\thepage}
\fancyfoot[RE,LO]{\sffamily\scriptsize\color{sbgray}Engineering Model \textbullet{} schematico rev.\ 1.45 \textbullet{} documento v1.1}
\fancypagestyle{plain}{\fancyhf{}\renewcommand{\headrulewidth}{0pt}%
  \fancyfoot[LE,RO]{\sffamily\small\bfseries\color{sbblue}\thepage}}
\renewcommand{\chaptermark}[1]{\markboth{\thechapter.\ #1}{}}

% ---------------------------------------------------------------- riquadri
\newtcolorbox{attenzione}[1][]{enhanced,breakable,colback=sbredbg,colframe=sbred,
  boxrule=0.9pt,arc=1.5pt,left=7pt,right=7pt,top=4pt,bottom=4pt,
  fonttitle=\sffamily\bfseries\small,coltitle=white,colbacktitle=sbred,
  title={Attenzione\ifx&#1&\else{} \textendash{} #1\fi}}
\newtcolorbox{regola}[1][]{enhanced,breakable,colback=sbgreenbg,colframe=sbgreen,
  boxrule=0.9pt,arc=1.5pt,left=7pt,right=7pt,top=4pt,bottom=4pt,
  fonttitle=\sffamily\bfseries\small,coltitle=white,colbacktitle=sbgreen,
  title={Regola\ifx&#1&\else{} \textendash{} #1\fi}}
\newtcolorbox{nota}[1][]{enhanced,breakable,colback=sbgraybg,colframe=sbgray,
  boxrule=0.6pt,arc=1.5pt,left=7pt,right=7pt,top=4pt,bottom=4pt,
  fonttitle=\sffamily\bfseries\small,coltitle=white,colbacktitle=sbgray,
  title={Nota\ifx&#1&\else{} \textendash{} #1\fi}}
\newtcolorbox{verifica}[1][]{enhanced,breakable,colback=sborangebg,colframe=sborange,
  boxrule=0.6pt,arc=1.5pt,left=7pt,right=7pt,top=4pt,bottom=4pt,
  fonttitle=\sffamily\bfseries\small,coltitle=white,colbacktitle=sborange,
  title={Da verificare\ifx&#1&\else{} \textendash{} #1\fi}}

% ---------------------------------------------------------------- macro
\newcommand{\net}[1]{\texttt{#1}}
\newcommand{\code}[1]{\texttt{#1}}
\newcommand{\fonte}[1]{\leavevmode{\footnotesize\color{sbgray}[#1]}}
\newcommand{\DV}{\textcolor{sborange}{\sffamily\bfseries\footnotesize da verificare}}
\newcommand{\Fid}[1]{\textbf{\textcolor{sbgreen}{#1}}}
\newcommand{\tabhead}{\rowcolor{sbblue}}
\newcommand{\hd}[1]{\textcolor{white}{\sffamily\bfseries\footnotesize #1}}
\newcolumntype{L}[1]{>{\raggedright\arraybackslash}p{#1}}
\newcolumntype{C}[1]{>{\centering\arraybackslash}p{#1}}
\newcolumntype{Y}{>{\raggedright\arraybackslash}X}
\setlength{\tabcolsep}{4pt}
\renewcommand{\arraystretch}{1.18}
\captionsetup{font={small,sf},labelfont={bf,color=sbblue},labelsep=period}
\setlist{itemsep=2pt,topsep=4pt}
\setlength{\parindent}{0pt}
\setlength{\parskip}{4pt plus 1pt}
\tikzset{timing/.style={x=3.2ex,y=2ex},
  blk/.style={draw=sbblue,thick,fill=sblight,rounded corners=2pt,align=center,
     font=\sffamily\footnotesize,minimum height=9mm,inner sep=3pt},
  blkr/.style={blk,draw=sbred,fill=sbredbg},
  blkg/.style={blk,draw=sbgreen,fill=sbgreenbg},
  blko/.style={blk,draw=sborange,fill=sborangebg}}
\tikzset{timing/rowdist=3.2ex}
% figure: niente pagine con una figura sola
\renewcommand{\topfraction}{0.9}\renewcommand{\bottomfraction}{0.8}
\renewcommand{\textfraction}{0.07}\renewcommand{\floatpagefraction}{0.75}
\setcounter{topnumber}{3}\setcounter{totalnumber}{4}
% riquadro "In parole semplici"
\newtcolorbox{semplice}[1][]{enhanced,breakable,colback=sblight,colframe=sbblue,
  boxrule=0.6pt,arc=1.5pt,left=7pt,right=7pt,top=4pt,bottom=4pt,
  fonttitle=\sffamily\bfseries\small,coltitle=white,colbacktitle=sbblue,
  title={In parole semplici\ifx&#1&\else{} \textendash{} #1\fi}}

\begin{document}
\pagenumbering{roman}
% ---------------------------------------------------------------- frontespizio
\begin{titlepage}
\newgeometry{margin=0mm}
\begin{tikzpicture}[remember picture,overlay]
  \fill[sbblue] (current page.north west) rectangle ([yshift=-95mm]current page.north east);
  \fill[sbblue2] ([yshift=-95mm]current page.north west) rectangle ([yshift=-99mm]current page.north east);
  \node[anchor=north west,text=white,font=\sffamily\bfseries\fontsize{40}{44}\selectfont]
     at ([xshift=25mm,yshift=-30mm]current page.north west) {SYMPHONY V2};
  \node[anchor=north west,text=white,font=\sffamily\Large]
     at ([xshift=25mm,yshift=-50mm]current page.north west) {Ricetrasmettitore UHF \SI{435}{\mega\hertz} per CubeSat};
  \node[anchor=north west,text=white!85!sbblue,font=\sffamily\large]
     at ([xshift=25mm,yshift=-62mm]current page.north west) {Engineering Model};
  % motivo grafico: catena RF stilizzata
  \begin{scope}[shift={([xshift=25mm,yshift=-150mm]current page.north west)},
     every node/.style={font=\sffamily\small}]
    \node[blk,minimum width=24mm] (a) at (0,0) {MCU\\STM32L562};
    \node[blk,minimum width=24mm] (b) at (34mm,0) {Si4464\\+ TCXO};
    \node[blkr,minimum width=24mm] (c) at (68mm,0) {PA\\QPA9510};
    \node[blk,minimum width=24mm] (d) at (102mm,0) {LPF\\FL201};
    \node[blko,minimum width=24mm] (e) at (136mm,0) {Switch\\HMC349A};
    \draw[-{Stealth},thick,sbblue] (a)--(b); \draw[-{Stealth},thick,sbblue] (b)--(c);
    \draw[-{Stealth},thick,sbblue] (c)--(d); \draw[-{Stealth},thick,sbblue] (d)--(e);
    \draw[thick,sbblue] (e.east)--++(6mm,0) node[right]{MMCX};
  \end{scope}
  \node[anchor=north west,font=\sffamily\bfseries\fontsize{30}{34}\selectfont,text=sbblue]
     at ([xshift=25mm,yshift=-185mm]current page.north west) {Manuale per il firmwarista};
  \node[anchor=north west,font=\sffamily\large,text=sbgray,text width=160mm]
     at ([xshift=25mm,yshift=-205mm]current page.north west)
     {Interfaccia hardware/firmware, mappa dei pin, sequenze RF, limiti di potenza e checklist di avvio};
  \node[anchor=south west,font=\sffamily\small,text=sbgray,text width=165mm,align=left]
     at ([xshift=25mm,yshift=22mm]current page.south west) {%
     \begin{tabular}{@{}ll@{}}
     \textbf{Versione del documento} & 1.0 del 04/11/2026\\
     \textbf{Riferimento hardware} & schematico KiCad 10, revisione \textbf{1.45} del log hardware (04/11/2026)\\
     \textbf{Netlist} & esportata con \code{kicad-cli 10.0} dallo schematico salvato il 04/11/2026 15:22\\
     \textbf{Stato} & Engineering Model \textendash{} con LED di debug montati\\
     \textbf{Lingua} & italiano
     \end{tabular}};
\end{tikzpicture}
\restoregeometry
\end{titlepage}

% ---------------------------------------------------------------- storico
\cleardoublepage
\begingroup\rowcolors{2}{white}{sbrow}
\thispagestyle{plain}
{\sffamily\Large\bfseries\color{sbblue} Storia del documento\par}
\vspace{4mm}
\begin{xltabular}{\linewidth}{@{}L{18mm}L{24mm}L{34mm}Y@{}}
\tabhead\hd{Versione}&\hd{Data}&\hd{Hardware}&\hd{Descrizione}\\
1.0 & 04/11/2026 & rev.\ 1.45 del log & Prima emissione. Ricavata dalla netlist dello schematico, dal log hardware (entry 000\ldots043), dalle revisioni indipendenti del 03/11 e del 04/11/2026 e dai datasheet dei componenti.\\
\bottomrule
\end{xltabular}\endgroup

\vspace{8mm}
{\sffamily\Large\bfseries\color{sbblue} Fonti usate\par}
\vspace{2mm}
\begin{itemize}
\item \textbf{[SCH]} Schematico KiCad 10 \code{SYMPHONY.kicad\_sch} con i fogli POWER, UHF, MCU, IF, CONNECTORS (= SENSORS), MOUNTING\_HOLES. Netlist \code{net.xml} e BOM \code{bom.csv} esportate nella cartella di questo manuale.
\item \textbf{[LOG]} Log hardware \code{report/simphony\_v2\_log} (entry 000\ldots043, fino alla rev.\ 1.45). Le note firmware F-01\ldots F-16 e le misure M-01\ldots M-06 vengono da qui.
\item \textbf{[REV-1103]}, \textbf{[REV-1104]} Revisioni indipendenti dello schematico del 03/11/2026 e del 04/11/2026 (\code{SYMPHONY\_V2\_review\_finale\_*.pdf}).
\item \textbf{[DS-xxx]} Datasheet dei componenti (elenco e pagine chiave in appendice~\ref{app:componenti}). Le pagine indicate sono le pagine del PDF.
\end{itemize}
Ogni valore numerico ha accanto la sua fonte tra parentesi quadre. Se un valore non è stato verificato sulla fonte è marcato \DV{} ed è raccolto nell'appendice~\ref{app:daverificare}.

\tableofcontents
\cleardoublepage
\pagenumbering{arabic}
\chapter{Introduzione}
\label{ch:intro}

\section{Scopo del documento}
Questo manuale descrive tutto quello che serve per scrivere il firmware della scheda SYMPHONY V2. La scheda è un ricetrasmettitore \textbf{solo UHF} a \SI{435}{\mega\hertz} per CubeSat. Non c'è una catena VHF.

Il manuale è scritto per chi non conosce il progetto. Non serve aprire lo schematico per scrivere il firmware: pin, livelli, tempi e limiti sono riportati qui con la loro fonte.

L'hardware descritto è l'\textbf{Engineering Model} (EM). I LED servono come riscontro visivo al banco. Sul \textbf{Flight Model} (FM) i LED di debug non si montano (capitolo~\ref{ch:led}).

\section{Come usare questo manuale}
\begin{enumerate}
\item Leggere il capitolo~\ref{ch:sistema} per avere il quadro della scheda.
\item Configurare clock e GPIO con la mappa dei pin del capitolo~\ref{ch:pin}. È la tabella di riferimento per CubeMX o per il codice di inizializzazione.
\item Scrivere i driver con il capitolo~\ref{ch:periferiche} (SPI, I\textsuperscript{2}C, UART, CAN, ADC, interrupt).
\item Scrivere il driver radio e la macchina a stati TX/RX con il capitolo~\ref{ch:rf}. È la parte più delicata: \textbf{un errore di sequenza può danneggiare il commutatore RF}.
\item Implementare watchdog, RTC e FRAM con il capitolo~\ref{ch:supervisione}.
\item Prima di consegnare il firmware, spuntare la checklist del capitolo~\ref{ch:notefw} e la tabella delle note F-01\ldots F-16.
\end{enumerate}

\subsection*{Convenzioni}
\begin{itemize}
\item \textbf{Nomi delle net.} Sono scritti così: \net{/MCU/UHF.EN\_TX\_ANT\_SW}. Sono i nomi esatti della netlist. La prima parte (\net{/MCU/}) è il foglio dello schematico in cui nasce il segnale; le net senza barra (\net{3V3}, \net{RF\_MR\_N}) sono globali, cioè hanno lo stesso nome su tutti i fogli.
\item \textbf{Sigle dei componenti.} U301, IC204, R243\ldots{} sono quelle stampate sullo schematico, sulla BOM (la distinta dei componenti) e sulla serigrafia della scheda.
\item \textbf{Livelli logici.} \textbf{H} = livello alto (\SI{3.3}{\volt}). \textbf{L} = livello basso (\SI{0}{\volt}). \textbf{Hi-Z} = alta impedenza: il pin non spinge né alto né basso, come se fosse staccato, e il livello lo decide una resistenza esterna (se c'è).
\item \textbf{Frecce nelle tabelle.} \enquote{A $\to$ B} vuol dire \enquote{prima A, poi B}. Per esempio \enquote{Hi-Z $\to$ H} = il pin è in alta impedenza dopo il reset e il firmware lo deve portare alto.
\item \textbf{Fonti.} Sono tra parentesi quadre, in grigio, per esempio \fonte{DS-Si4464 p.\,17}. Le sigle:
 \begin{itemize}
 \item \textbf{DS-xxx} = datasheet del componente xxx, \enquote{p.} = pagina del PDF (elenco nell'appendice~\ref{app:componenti});
 \item \textbf{SCH} = schematico KiCad del progetto;
 \item \textbf{LOG} = log delle modifiche hardware; \enquote{rev.\ 1.45} = revisione del log, \enquote{§41} = paragrafo 41 del log;
 \item \textbf{REV-1103}, \textbf{REV-1104} = revisioni indipendenti dello schematico del 03/11 e del 04/11/2026.
 \end{itemize}
\item \textbf{Codici del log.} Nel log ogni punto ha un codice; qui si usano così:
 \begin{itemize}
 \item \textbf{F-01\ldots F-16} = note per il firmware, cioè cose che il firmware \emph{deve} fare (elenco completo nella tabella~\ref{tab:notefw});
 \item \textbf{M-01\ldots M-06} = misure da fare sul prototipo (tabella~\ref{tab:collaudo});
 \item \textbf{R-xx}, \textbf{S-xx}, \textbf{P-xx}, \textbf{D1}, numeri come \enquote{1.15-01} = punti del log o delle revisioni da cui viene un'informazione. Servono solo per ritrovare la fonte: non serve conoscerli per scrivere il firmware;
 \item \textbf{I-x} e \textbf{V-x} = incongruenze e dati da verificare raccolti da questo manuale (appendice~\ref{app:daverificare}).
 \end{itemize}
\item \DV{} indica un dato che non è stato possibile controllare sulla fonte. Tutti questi punti sono raccolti nell'appendice~\ref{app:daverificare}.
\item \textbf{Modo A / modo B}: i due modi in cui si può montare il reset della radio (paragrafo~\ref{sec:sdn}). Di serie la scheda è in modo A.
\end{itemize}

\subsection*{I cinque tipi di riquadro}
\begin{semplice}[esempio]
Spiegazione a parole, senza dettagli tecnici, di quello che segue.
\end{semplice}
\begin{attenzione}[esempio]
Rischio di danno all'hardware o di malfunzionamento in volo. Va rispettato sempre.
\end{attenzione}
\begin{regola}[esempio]
Requisito che il firmware deve implementare. Di solito corrisponde a una nota F-xx del log.
\end{regola}
\begin{nota}[esempio]
Spiegazione o consiglio. Non è vincolante.
\end{nota}
\begin{verifica}[esempio]
Dato non confermato o incongruenza trovata tra le fonti.
\end{verifica}

\section{Glossario}
Le sigle sono in ordine alfabetico. Tra parentesi c'è la sigla del componente sulla scheda.

\begin{xltabular}{\linewidth}{@{}L{26mm}Y@{}}
\tabhead\hd{Termine}&\hd{Significato}\\
\endhead
abs.\ max & \emph{Absolute maximum}: valore che non va mai superato, nemmeno per un istante, pena il danno del componente.\\
ADC &Convertitore analogico-digitale dell'MCU. Qui legge le due temperature e la tensione del 3V3.\\
AF & \emph{Alternate Function}: un pin dell'STM32 può essere un GPIO o una periferica (SPI, UART\ldots). Il numero AF (AF0\ldots AF15) dice quale periferica, e si scrive nei registri \code{GPIOx\_AFRL/AFRH}.\\
APC & \emph{Automatic Power Control}: pin di controllo del PA QPA9510 (APC1/APC2). A \SI{0}{\volt} il PA è a riposo, a circa \SI{2.8}{\volt} lavora alla potenza massima. Su questa scheda l'APC è solo acceso/spento, non regolabile.\\
BOM & \emph{Bill Of Materials}: elenco dei componenti da montare.\\
CS, NSEL & \emph{Chip Select}: pin che seleziona un dispositivo SPI. Basso = dispositivo selezionato. NSEL è il nome del CS sul Si4464.\\
CTS & \emph{Clear To Send} del Si4464: byte \code{0xFF} che dice \enquote{ho finito il comando precedente, puoi mandarne un altro}.\\
dBm & Potenza riferita a \SI{1}{\milli\watt}: \SI{0}{\dBm} = \SI{1}{\milli\watt}, \SI{10}{\dBm} = \SI{10}{\milli\watt}, \SI{30}{\dBm} = \SI{1}{\watt}. Ogni \SI{3}{dB} in più la potenza raddoppia.\\
DNP & \emph{Do Not Populate}: componente disegnato sullo schematico ma non montato.\\
EM / FM & \emph{Engineering Model} (prototipo da banco, con LED) / \emph{Flight Model} (scheda di volo).\\
EXTI & Controllore degli interrupt esterni dell'STM32: genera un interrupt quando un pin cambia livello.\\
FRAM & Memoria non volatile FM25V10 (U401), \SI{1}{Mbit}. Come una EEPROM, ma si scrive subito e quasi senza usura.\\
GPIO & Pin usato come semplice ingresso o uscita digitale.\\
HSE & Oscillatore esterno ad alta velocità dell'MCU (quarzo Y1, \SI{24}{\mega\hertz}).\\
Hi-Z & Alta impedenza: pin né alto né basso, come scollegato.\\
Hot switching & Commutare lo switch RF mentre c'è potenza RF che lo attraversa. Lo stressa: va evitato.\\
I2C & Bus seriale a due fili (SCL clock, SDA dati) con resistenze di pull-up.\\
LNA & \emph{Low Noise Amplifier}: amplificatore di ricezione TQP3M9036 (Q201).\\
LPF & Filtro passa-basso FL201 (Mini-Circuits LFCG-490+) dopo il PA: toglie le armoniche.\\
MMCX & Connettore d'antenna J201.\\
MR & \emph{Manual Reset}: ingresso \(\overline{\mathrm{MR}}\) del supervisore TPS3825 (IC205). La barra sopra il nome vuol dire \enquote{attivo basso}.\\
nIRQ & Uscita di interrupt del Si4464, attiva bassa (la \enquote{n} iniziale vuol dire attivo basso).\\
Open-drain (OD) & Uscita che può solo tirare la linea verso massa; il livello alto lo dà una resistenza di pull-up.\\
PA & \emph{Power Amplifier}: amplificatore di potenza QPA9510 (IC2).\\
PA\_PWR\_LVL & Impostazione del Si4464 che regola la sua potenza d'uscita (128 passi). Il PA amplifica quello che esce dal Si4464.\\
POR & \emph{Power-On Reset}: reset automatico all'accensione.\\
PT1000 & Sensore di temperatura a resistenza: \SI{1000}{\ohm} a \SI{0}{\celsius}, cresce di circa \SI{3.9}{\ohm} per grado.\\
Pull-up / pull-down & Resistenza che porta una linea alta (verso 3V3) o bassa (verso massa) quando nessuno la pilota.\\
Push-pull & Uscita che spinge attivamente sia alto sia basso.\\
RTC & \emph{Real Time Clock}: orologio DS3231M (U404).\\
SDN & Pin di shutdown del Si4464: alto = radio spenta e configurazione persa.\\
SPI & Bus seriale sincrono a 4 fili: SCK clock, MOSI dati verso il dispositivo, MISO dati verso l'MCU, CS selezione.\\
SWD & \emph{Serial Wire Debug}: interfaccia del debugger (programmazione e breakpoint).\\
Switch & Commutatore RF HMC349A (IC202): collega l'antenna al ramo TX o al ramo RX.\\
TCXO & Oscillatore al quarzo compensato in temperatura Y201, \SI{32}{\mega\hertz}: è il clock del Si4464.\\
VREFINT & Tensione di riferimento interna dell'STM32. Leggendola con l'ADC si ricava la tensione del 3V3.\\
Watchdog & Circuito che resetta l'MCU se il firmware smette di dare segni di vita. Qui è esterno: STWD100 (U403).\\
WDI / WDO & Ingresso degli impulsi / uscita di reset del watchdog STWD100 (U403).\\
WDS & \emph{Wireless Development Suite} di Silicon Labs: programma che genera la configurazione del Si4464 (\code{radio\_config.h}).\\
\bottomrule
\end{xltabular}

\chapter{Panoramica del sistema}
\label{ch:sistema}

\begin{semplice}
SYMPHONY V2 è una radio UHF per un CubeSat. Il microcontrollore (MCU) STM32L562 comanda la radio Si4464, che trasmette e riceve a \SI{435}{\mega\hertz}. In trasmissione il segnale viene amplificato da un amplificatore di potenza (PA). In ricezione passa da un amplificatore a basso rumore (LNA). Un commutatore RF (switch) collega l'antenna a uno dei due rami. Intorno ci sono una memoria (FRAM), un orologio (RTC), un watchdog, due sensori di temperatura e le interfacce verso il satellite (CAN, I2C) e verso il PC (USB).
\end{semplice}

\section{Schema a blocchi}
Lo schema è diviso in tre figure, per non affollare il disegno:
\begin{itemize}
\item figura~\ref{fig:blocchi}: l'MCU e tutto quello che le è collegato;
\item figura~\ref{fig:catena-rf}: il percorso del segnale radio;
\item figura~\ref{fig:alim-blocchi}: da dove arriva l'alimentazione.
\end{itemize}
Sulle frecce è scritta \textbf{la periferica dell'MCU} (per esempio FDCAN1), non il numero del pin. I pin di ogni periferica sono nella mappa dei pin (capitolo~\ref{ch:pin}) e nelle tabelle del capitolo~\ref{ch:periferiche}. Tutti i collegamenti sono presi dalla netlist dello schematico \fonte{SCH}.

\begin{figure}[htbp]
\centering
\begin{tikzpicture}[x=1mm,y=1mm,
   per/.style={blk,minimum width=34mm,minimum height=11mm},
   lab/.style={font=\sffamily\scriptsize,inner sep=1.5pt},
   bus/.style={{Stealth[length=2mm]}-{Stealth[length=2mm]},thick,sbblue2},
   uni/.style={-{Stealth[length=2mm]},thick,sbblue2},
   ctl/.style={-{Stealth[length=2mm]},thick,sbgreen}]
% MCU
\node[blk,minimum width=36mm,minimum height=104mm,font=\sffamily\small] (mcu) at (0,0)
  {\textbf{U301}\\\textbf{STM32L562}\\[3pt]microcontrollore\\[2pt]\footnotesize quarzo \SI{24}{\mega\hertz}};
% ---- sinistra
\node[per] (fram) at (-62,40) {U401 FM25V10\\memoria FRAM};
\node[per] (rtc)  at (-62,20) {U404 DS3231M\\orologio (RTC)};
\node[per] (wdt)  at (-62,0)  {U403 STWD100\\watchdog esterno};
\node[per] (adc)  at (-62,-20){IC1, IC5 TLV9051\\2 sensori PT1000};
\node[per] (swd)  at (-62,-40){J1 TC2030\\debugger};
\draw[bus] (fram.east) -- node[lab,above]{SPI2} (fram.east-|mcu.west);
\draw[bus] ($(rtc.east)+(0,2)$) -- node[lab,above]{I2C2} ($(rtc.east-|mcu.west)+(0,2)$);
\draw[uni] ($(rtc.east)+(0,-2)$) -- node[lab,below]{GPIO: avviso} ($(rtc.east-|mcu.west)+(0,-2)$);
\draw[uni] ($(wdt.east-|mcu.west)+(0,2)$) -- node[lab,above]{GPIO: impulso} ($(wdt.east)+(0,2)$);
\draw[ctl] ($(wdt.east)+(0,-2)$) -- node[lab,below,text=sbgreen]{reset (NRST)} ($(wdt.east-|mcu.west)+(0,-2)$);
\draw[uni] (adc.east) -- node[lab,above]{ADC1} (adc.east-|mcu.west);
\draw[bus] (swd.east) -- node[lab,above]{SWD} (swd.east-|mcu.west);
% ---- destra
\node[per] (si)  at (62,45)  {IC204 Si4464\\radio};
\node[per,blko] (sup) at (62,27)  {IC205 TPS3825\\supervisore radio};
\node[per,blko] (cmd) at (62,9)   {Q701 + U701\\switch, PA, LNA};
\node[per] (ch)  at (62,-9)  {U402 CH340E\\verso USB-C (J15)};
\node[per] (can) at (62,-27) {IC501 TCAN3414\\verso J202 e dock};
\node[per] (i2c) at (62,-45) {bus I2C esterno\\J203 e dock J708};
\draw[bus] ($(si.west-|mcu.east)+(0,2)$) -- node[lab,above]{SPI1} ($(si.west)+(0,2)$);
\draw[uni] ($(si.west)+(0,-2)$) -- node[lab,below]{EXTI (interrupt)} ($(si.west-|mcu.east)+(0,-2)$);
\draw[uni] (sup.west-|mcu.east) -- node[lab,above]{GPIO: reset radio} (sup.west);
\draw[uni] (cmd.west-|mcu.east) -- node[lab,above]{GPIO: TX/RX} (cmd.west);
\draw[bus] (ch.west-|mcu.east) -- node[lab,above]{USART1} (ch.west);
\draw[bus] (can.west-|mcu.east) -- node[lab,above]{FDCAN1} (can.west);
\draw[bus] (i2c.west-|mcu.east) -- node[lab,above]{I2C1} (i2c.west);
\draw[ctl] (sup.north) -- node[lab,right,text=sbgreen]{SDN} (si.south);
\end{tikzpicture}
\caption{L'MCU e le sue periferiche. Sulle frecce c'è la periferica dell'MCU usata. Blu: dati (la punta indica il verso, due punte = nei due versi). Verde: segnale che agisce da solo, senza passare dal firmware. Arancio: blocchi che proteggono la radio.}
\label{fig:blocchi}
\end{figure}

\begin{nota}[come leggere la figura~\ref{fig:blocchi}]
\begin{itemize}
\item \textbf{GPIO}: pin usato come semplice ingresso o uscita digitale, non come periferica. Dopo i due punti c'è a che cosa serve.
\item \textbf{GPIO: avviso}: il DS3231M abbassa questo pin quando la sua alimentazione cade (paragrafo~\ref{sec:rtc-rst}).
\item \textbf{GPIO: impulso}: il firmware deve cambiarne il livello almeno ogni secondo, altrimenti il watchdog resetta l'MCU tramite \textbf{NRST} (capitolo~\ref{ch:supervisione}).
\item \textbf{EXTI}: linea di interrupt esterno. Il Si4464 abbassa il pin quando ha un evento da segnalare (pacchetto inviato, pacchetto ricevuto\ldots).
\item \textbf{SDN}: il supervisore IC205 spegne la radio se il 3V3 scende troppo o se il firmware lo chiede con \enquote{GPIO: reset radio} (paragrafo~\ref{sec:sdn}).
\item \textbf{GPIO: TX/RX}: un solo pin decide se la scheda trasmette o riceve. Comanda insieme switch, PA e LNA (figura~\ref{fig:catena-rf} e paragrafo~\ref{sec:rf-arch}).
\end{itemize}
\end{nota}

\begin{figure}[htbp]
\centering
\begin{tikzpicture}[x=1mm,y=1mm,
   rb/.style={blk,minimum width=22mm,minimum height=11mm},
   rf/.style={-{Stealth[length=2mm]},very thick,sborange},
   lab/.style={font=\sffamily\scriptsize,inner sep=1.5pt},
   ctl/.style={-{Stealth[length=2mm]},thick,sbgreen},
   cl/.style={font=\sffamily\scriptsize,text=sbgreen,align=center}]
\node[blk,minimum width=20mm,minimum height=34mm] (si) at (0,-9) {IC204\\Si4464\\radio};
\node[blk,minimum width=20mm] (tcxo) at (0,22) {Y201 TCXO\\\SI{32}{\mega\hertz}};
\draw[-{Stealth[length=2mm]},thick,sbblue2] (tcxo) -- node[lab,right]{clock} (si);
% TX
\node[rb] (att) at (31,0)  {attenuatore\\\SI{10.8}{dB}};
\node[rb,blkr] (pa) at (58,0)  {IC2 QPA9510\\PA};
\node[rb] (lpf) at (85,0)  {FL201\\passa-basso};
% RX
\node[rb] (mat) at (31,-18) {adattamento\\RX};
\node[rb] (lna) at (58,-18) {Q201\\LNA};
\node[rb] (frx) at (85,-18) {filtro RX\\passa-basso};
\node[blko,minimum width=20mm,minimum height=34mm] (sw) at (114,-9) {IC202\\HMC349A\\switch};
\node[draw=sbgray,thick,circle,inner sep=1.5pt,font=\sffamily\scriptsize] (ant) at (139,-9) {MMCX};
\draw[rf] (si.east|-att) -- node[lab,above]{TX} (att);
\draw[rf] (att) -- (pa); \draw[rf] (pa) -- (lpf);
\draw[rf] (lpf) -- node[lab,above]{RF2} (lpf-|sw.west);
\draw[rf] (frx-|sw.west) -- node[lab,above]{RF1} (frx);
\draw[rf] (frx) -- (lna); \draw[rf] (lna) -- (mat);
\draw[rf] (mat) -- node[lab,above]{RX} (mat-|si.east);
\draw[very thick,sborange,{Stealth[length=2mm]}-{Stealth[length=2mm]}] (sw) -- node[lab,above]{RFC} (ant);
% comandi
\node[cl] (cpa) at (58,16) {comando TX/RX\\(accende il PA, tramite Q701)};
\node[cl] (csw) at (114,16) {comando TX/RX\\(sceglie il ramo)};
\node[cl] (cln) at (58,-34) {comando TX/RX\\(spegne l'LNA, tramite U701)};
\draw[ctl] (cpa) -- (pa); \draw[ctl] (csw) -- (sw);
\draw[ctl] (cln) -- (lna);
\end{tikzpicture}
\caption{Percorso del segnale radio. Sopra la trasmissione (TX), sotto la ricezione (RX). Lo switch collega l'antenna (connettore MMCX) a un ramo o all'altro: RF2 = trasmissione, RF1 = ricezione. In verde il comando TX/RX: è lo stesso pin dell'MCU per tutti e tre i blocchi.}
\label{fig:catena-rf}
\end{figure}

\begin{figure}[htbp]
\centering
\begin{tikzpicture}[x=1mm,y=1mm,
   pb/.style={blk,minimum width=30mm,minimum height=11mm},
   pwr/.style={-{Stealth[length=2mm]},very thick,sbred!80},
   lab/.style={font=\sffamily\scriptsize,inner sep=1.5pt},
   ctl/.style={-{Stealth[length=2mm]},thick,sbgreen}]
\node[pb] (dock) at (0,10)  {J708 dock\\(dal satellite)};
\node[pb] (usb)  at (0,-10) {J15 USB-C\\(dal PC)};
\node[pb] (mux)  at (46,0) {IC4 TPS2121\\scelta della sorgente};
\node[pb] (buck) at (94,0) {IC3 TPSM863257\\convertitore \SI{3.3}{\volt}};
\node[font=\sffamily\scriptsize,align=left,anchor=west] (load) at (120,0) {MCU, radio, PA,\\memorie, sensori};
\node[font=\sffamily\scriptsize,align=left,anchor=west] (swv) at (120,-16) {alimentazione\\dello switch RF};
\draw[pwr] (dock.east) -| ($(mux.west)+(-6,3)$) -- ($(mux.west)+(0,3)$);
\node[lab,above] at ($(dock.east)+(8,0)$) {5\,V, priorità};
\draw[pwr] (usb.east) -| ($(mux.west)+(-6,-3)$) -- ($(mux.west)+(0,-3)$);
\node[lab,below] at ($(usb.east)+(8,0)$) {5\,V USB};
\draw[pwr] (mux) -- node[lab,above]{5V0} (buck);
\draw[pwr] (buck) -- node[lab,above]{3V3} (load);
\draw[pwr] (69,0) |- node[lab,above,pos=0.75]{5V0} (swv);
\fill[sbred!80] (69,0) circle (0.8);
\draw[ctl] (dock.north) |- (94,21) node[lab,above,pos=0.75,text=sbgreen]{EN\_DOCK: accende il convertitore} -- (buck.north);
\end{tikzpicture}
\caption{Alimentazione. Il dock ha la precedenza sulla USB-C. Il convertitore del 3V3 si accende \textbf{solo} se il dock alza il segnale EN\_DOCK: con la sola USB-C l'MCU resta spenta (capitolo~\ref{ch:power}).}
\label{fig:alim-blocchi}
\end{figure}

\section{Componenti principali}
La colonna \enquote{Rif.} è la sigla del componente sullo schematico e sulla scheda.

\begin{xltabular}{\linewidth}{@{}L{17mm}L{37mm}Y@{}}
\tabhead\hd{Rif.}&\hd{Parte}&\hd{A che cosa serve, per il firmware}\\
\endhead
U301 & STM32L562RET6P & Il microcontrollore (MCU), Cortex-M33, contenitore LQFP64 da 64 pin. La lettera \enquote{P} finale indica la versione predisposta per un alimentatore esterno (SMPS): qui non è usato e i pin VDD12 non sono collegati \fonte{REV-1104 §2.3}.\\
IC204 & Si4464-B1B-FMR & La radio (transceiver sub-GHz): modula, demodula e gestisce i pacchetti. Si comanda via SPI1. Il clock arriva dal TCXO a \SI{32}{\mega\hertz}.\\
Y201 & ECS-TXO-20CSMV4-320-AY & TCXO = oscillatore compensato in temperatura, \SI{32}{\mega\hertz}. Precisione \SI{\pm0.5}{\ppm}, garantito da \SIrange{-30}{85}{\celsius} \fonte{DS-ECS p.\,1}.\\
IC2 & Qorvo QPA9510 & PA = amplificatore di potenza della trasmissione. Si accende quando il transistor Q701 alimenta il suo pin di controllo APC.\\
IC202 & ADI HMC349A & Switch RF: collega l'antenna al ramo TX (porta RF2) o al ramo RX (porta RF1).\\
Q201 & Qorvo TQP3M9036 & LNA = amplificatore a basso rumore della ricezione. Ha un pin per spegnerlo.\\
U701 & 74AUP2G14 & Porta logica invertente: da PC2 ricava il comando opposto per l'LNA (net \net{EN\_RX\_LNA}).\\
IC205 & TI TPS3825-33 & Supervisore: tiene la radio spenta se il 3V3 è troppo basso o se il firmware porta basso PC4.\\
IC4 & TI TPS2121 & Sceglie l'alimentazione a \SI{5}{\volt}: il dock (ingresso IN1) ha la precedenza sulla USB-C (ingresso IN2).\\
IC3 & TI TPSM863257 & Convertitore (buck) da 5V0 a 3V3. Si accende solo con \net{EN\_DOCK} alto.\\
U403 & ST STWD100NYWY3F & Watchdog esterno: resetta l'MCU se il firmware non manda impulsi. Tempo tipico \SI{1.6}{\second}, minimo \SI{1.12}{\second}.\\
U404 & Maxim DS3231M & Orologio (RTC) su I2C2. Mantiene l'ora per qualche ora a scheda spenta grazie a dei supercondensatori.\\
U401 & Infineon FM25V10 & Memoria non volatile FRAM da \SI{1}{Mbit} (\SI{128}{\kilo\byte}) su SPI2. Si scrive senza attese.\\
IC501 & TI TCAN3414 & Transceiver CAN: converte FDCAN1 nei livelli del bus CAN.\\
U402 & WCH CH340E & Convertitore USB$\leftrightarrow$seriale su USART1: console e caricamento del firmware.\\
IC1, IC5 & TI TLV9051 & Amplificatori operazionali delle due sonde di temperatura PT1000 (PA e buck).\\
\bottomrule
\end{xltabular}

\section{Fatti essenziali da ricordare}
\begin{itemize}
\item La radio lavora a \SI{435}{\mega\hertz}, dentro la banda \SIrange{353}{479}{\mega\hertz} del Si4464 \fonte{DS-Si4464 p.\,9}.
\item Un solo pin, \textbf{PC2}, comanda insieme switch RF, PA e LNA \fonte{LOG §41, R-04}. PC2 basso = trasmissione, PC2 alto = ricezione.
\item Il 3V3 esiste solo se il dock alza \net{EN\_DOCK}. Con la sola USB-C l'MCU resta spenta (capitolo~\ref{ch:power}).
\item Il watchdog esterno è \textbf{sempre attivo} in volo: il primo impulso su PC6 va dato entro \SI{1.12}{\second} dall'accensione \fonte{LOG F-12}.
\item Nessun LED è comandato dall'MCU. D203 mostra se la radio è fuori reset, D9 la presenza del 3V3 (solo con il ponticello JP1 chiuso).
\end{itemize}

\chapter{Alimentazioni e accensione}
\label{ch:power}

\begin{semplice}
La scheda prende i \SI{5}{\volt} dal dock del satellite oppure dalla USB-C; se ci sono tutti e due vince il dock. Dai \SI{5}{\volt} un convertitore ricava i \SI{3.3}{\volt} (il \enquote{3V3}) che alimentano quasi tutto. Il convertitore però si accende solo se il dock gli dà il permesso con il segnale EN\_DOCK: con la sola USB-C la scheda resta spenta.
\end{semplice}

\section{Rail di alimentazione}
Un \enquote{rail} è una linea di alimentazione. Nella tabella: \textbf{Rail} = nome della net; \textbf{Generato da} = il componente o il connettore che la fornisce; \textbf{Valore} = tensione nominale ed eventuale intervallo garantito; \textbf{Carichi principali} = cosa alimenta.

\begin{xltabular}{\linewidth}{@{}L{24mm}L{30mm}L{30mm}Y@{}}
\tabhead\hd{Rail}&\hd{Generato da}&\hd{Valore}&\hd{Carichi principali}\\
\endhead
\net{5V} & dock J708 pin 1 e 3 & \SI{5}{\volt} nominali & ingresso IN1 di IC4 (priorità)\\
\net{/POWER/VBUS} & USB-C J15 & \SI{5}{\volt}. Le resistenze da \SI{5.1}{\kilo\ohm} sui pin CC1/CC2 dicono al PC che la scheda è un dispositivo USB 2.0 & ingresso IN2 di IC4\\
\net{5V0} & IC4 TPS2121 (OUT) & uguale all'ingresso scelto & IC3 (buck), VDD di IC202 (switch RF)\\
\net{3V3} & IC3 TPSM863257 & \SI{3.30}{\volt} nominali, \SIrange{3.25}{3.405}{\volt} \fonte{REV-1104 §2.1} & tutta la logica, Si4464, TCXO, PA, LNA, sensori\\
\net{RTC\_VBAT} & 3V3 via D401 (BAT54WS) e R412 \SI{1}{\kilo\ohm} & $\approx$\SI{3.15}{\volt} \fonte{SCH nota MCU} & VBAT di U404, 10 supercondensatori CPH3225A (\SI{110}{\milli\farad})\\
\bottomrule
\end{xltabular}

Il pin VBAT dell'MCU (pin 1) e VDDUSB (pin 48) sono collegati direttamente al 3V3 \fonte{SCH}. Il dominio di backup dell'STM32 non ha una batteria propria: l'RTC interno dell'MCU perde l'ora quando il 3V3 si spegne. Il riferimento di tempo è il DS3231M (capitolo~\ref{ch:supervisione}).

I pin \net{+12V}, \net{+3V3\_CONN}, \net{RS485.A1/B1} del dock non sono usati su questa scheda \fonte{LOG R-12; SCH nota IF}.

\section{Power mux TPS2121: priorità al dock}
\begin{itemize}
\item \textbf{Priorità IN1 (dock).} PR1 = $5 \times 4{,}7/(10+4{,}7) = \SI{1.60}{\volt}$, sopra $V_\text{REF} = \SI{1.06}{\volt}$: con il dock presente vince sempre il dock \fonte{LOG D1; REV-1104 §2.1}. Il dock perde la priorità quando scende sotto $1{,}06 \times 14{,}7/4{,}7 = \SI{3.32}{\volt}$.
\item \textbf{Sovratensione IN1.} Soglia OV1 = $1{,}06 \times (22+4{,}7)/4{,}7 = \SI{6.02}{\volt}$ nominali \fonte{LOG rev.\ 1.44, P-08}. OV2 (USB) è disattivata (a GND).
\item \textbf{Limite di corrente.} ILIM con R44 = \SI{33}{\kilo\ohm}: \SI{3.21}{\ampere} \fonte{REV-1104 §2.1}.
\item \textbf{Stato.} Il pin ST è a GND: \textbf{il firmware non sa quale ingresso è attivo}.
\end{itemize}

\section{Buck TPSM863257 e segnale EN\_DOCK}
L'ingresso EN di IC3 è pilotato da \net{/IF/EN\_DOCK} (J708 pin 15) attraverso R46 \SI{10}{\kilo\ohm}. EN ha un pull-down interno da \SI{2}{\mega\ohm}: \textbf{senza EN\_DOCK alto il 3V3 resta spento} \fonte{LOG Parte B; REV-1104 §2.1}. È una scelta di progetto confermata \fonte{LOG D1}.

Il soft-start interno dura \SI{1.6}{\milli\second} tipici \fonte{DS-TPSM863257 p.\,7}. L'uscita PG non è collegata.

\begin{attenzione}[alimentazione al banco]
Con la sola USB-C il 5V0 c'è, ma il 3V3 no: l'MCU, il CH340E e la console sono spenti. Al banco bisogna fornire anche EN\_DOCK (EN max raccomandato \SI{5.5}{\volt}) \fonte{REV-1104 §2.1}. In trasmissione il 5V0 assorbe circa \SI{1}{\ampere}: una porta USB 2.0 di default dà \SI{500}{\milli\ampere}. Usare il dock o un alimentatore da almeno \SI{1.5}{\ampere} \fonte{REV-1104 P22}.
\end{attenzione}

\section{Che cosa vede il firmware}
L'hardware non porta all'MCU nessun segnale di stato dell'alimentazione (né ST del TPS2121, né PG del buck, né una misura del 5\,V). Il firmware ha a disposizione:
\begin{itemize}
\item \textbf{VDDA tramite VREFINT} (ADC1 canale 0):
\[ V_\text{DDA} = \SI{3.0}{\volt}\times\frac{\text{VREFINT\_CAL}}{\text{VREFINT\_DATA}} \]
VREFINT\_CAL si legge all'indirizzo \code{0x0BFA05AA}; è misurato in fabbrica a \SI{3.0}{\volt} e \SI{30}{\celsius} \fonte{DS-STM32L562 p.\,65}. Il 3V3 è anche VDDA, quindi questa è la misura del rail 3V3.
\item \textbf{Causa del reset} in \code{RCC\_CSR} (BOR, pin NRST, IWDG\ldots) \fonte{LOG F-13}.
\item \textbf{Pin RST del DS3231M} su PB9: va basso quando il VCC dell'RTC scende sotto la soglia di power-fail $V_\text{PF}$ \fonte{DS-DS3231M p.\,7}. Il valore numerico di $V_\text{PF}$ non è stato estratto dal datasheet: \DV.
\item \textbf{Temperatura del buck} (PT1000 su PC1) e del PA (PT1000 su PB1).
\item \textbf{Risposta del Si4464} dopo ogni TX: se ha perso la configurazione c'è stato un calo del 3V3 (nota F-10, paragrafo~\ref{sec:rf-supervisore}).
\end{itemize}

\section{Sequenza di accensione}
\label{sec:powerup}
La figura~\ref{fig:powerup} riassume cosa succede dal momento in cui il dock alza EN\_DOCK. I tempi sono tipici; i valori minimi e massimi sono nella tabella~\ref{tab:powerup}.

\begin{figure}[!htb]
\centering
\begin{tikzpicture}[x=0.72mm,y=1mm,font=\sffamily\scriptsize,
   sig/.style={thick},dim/.style={{Stealth[length=1.6mm]}-{Stealth[length=1.6mm]},sbgray}]
% righe: base (livello basso) a y = -13*i, livello alto a +5
\foreach \i/\name/\sub in {0/EN\_DOCK/dal dock,1/3V3/uscita del buck,2/NRST/reset MCU,3/WDI/impulsi del firmware,4/SDN/alto = radio spenta (modo A),5/RST DS3231M/avviso dell'RTC,6/WDO/reset dal watchdog}{
   \node[anchor=east,align=right] at (-2,-13*\i+2.5) {\textbf{\name}\\[-1pt]{\color{sbgray}\sub}};
   \draw[sbgray!30] (0,-13*\i) -- (150,-13*\i);}
% linee verticali degli eventi
\foreach \x/\t in {8/{A},14/{B},30/{C},72/{D}}{
   \draw[sbgray!70,densely dashed] (\x,8) -- (\x,-82);
   \node[circle,draw=sbgray,inner sep=0.8pt,fill=white] at (\x,10) {\t};}
% EN_DOCK
\draw[sig,sbgreen] (0,0) -- (8,0) -- (8,5) -- (150,5);
% 3V3
\draw[sig,sbred] (0,-13) -- (10,-13) -- (14,-8) -- (150,-8);
\draw[dim] (8,-15.5) -- (14,-15.5); \node[anchor=west,fill=white,inner sep=1pt] at (15,-15.5) {soft-start \SI{1.6}{\milli\second}};
% NRST
\draw[sig,sbblue] (0,-26) -- (14,-26) -- (14,-21) -- (150,-21);
% WDI
\draw[sig,sbblue] (0,-39) -- (30,-39) -- (30,-34) -- (60,-34) -- (60,-39) -- (90,-39) -- (90,-34) -- (120,-34) -- (120,-39) -- (150,-39);
\draw[dim] (14,-41.5) -- (30,-41.5); \node[anchor=west,fill=white,inner sep=1pt] at (31,-41.5) {primo fronte entro \SI{1.12}{\second}};
% SDN
\draw[sig,sborange] (0,-52) -- (14,-52) -- (14,-47) -- (72,-47) -- (72,-52) -- (150,-52);
\draw[dim] (14,-54.5) -- (72,-54.5); \node[anchor=west,fill=white,inner sep=1pt] at (73,-54.5) {$t_d$ = \SIrange{120}{300}{\milli\second}};
% RST DS3231M
\draw[sig,sbgreen] (0,-65) -- (78,-65) -- (78,-60) -- (150,-60);
\draw[dim] (14,-67.5) -- (78,-67.5); \node[anchor=west,fill=white,inner sep=1pt] at (79,-67.5) {$t_\text{REC}$ fino a \SI{300}{\milli\second}};
% WDO
\draw[sig,sbgreen] (0,-78) -- (14,-78) -- (14,-73) -- (150,-73);
\node[anchor=north,inner sep=1pt] at (118,-73.4) {resta alto: il firmware manda gli impulsi};
\end{tikzpicture}
\caption{Sequenza di accensione (schema non in scala). A: il dock alza EN\_DOCK. B: il 3V3 è pronto e l'MCU parte. C: primo impulso del firmware al watchdog. D: il supervisore libera la radio. Dopo C il firmware deve cambiare il livello di WDI almeno una volta al secondo. Linea bassa = livello logico 0, linea alta = livello logico 1.}
\label{fig:powerup}
\end{figure}

\begin{xltabular}{\linewidth}{@{}L{8mm}L{48mm}L{36mm}Y@{}}
\caption{Eventi all'accensione e tempi con fonte.}\label{tab:powerup}\\
\tabhead\hd{\#}&\hd{Evento}&\hd{Tempo}&\hd{Fonte e conseguenza per il firmware}\\
\endfirsthead
\tabhead\hd{\#}&\hd{Evento}&\hd{Tempo}&\hd{Fonte e conseguenza per il firmware}\\
\endhead
1 & Soft-start del buck, 3V3 sale & \SI{1.6}{\milli\second} tipici & \fonte{DS-TPSM863257 p.\,7}\\
2 & POR/BOR dell'MCU, boot da flash (BOOT0 = L tramite R3 \SI{10}{\kilo\ohm} a GND) & \textendash & \fonte{SCH}\\
3 & IC205 tiene SDN del Si4464 alto dopo che il 3V3 supera la soglia & $t_d$ = \SI{120}{}/\SI{200}{}/\SI{300}{\milli\second} (min/tip/max) & \fonte{DS-TPS3825 p.\,8}. Non usare il Si4464 prima di $t_d$ max: vedi F-05.\\
4 & Soglia del supervisore IC205 & \SI{2.88}{}/\SI{2.93}{}/\SI{3.00}{\volt} & \fonte{LOG §46; DS-TPS3825 p.\,6}\\
5 & Il Si4464 esce dallo shutdown: POR e calibrazioni interne, poi serve POWER\_UP & Shutdown$\to$TX o RX: \SI{15}{\milli\second} (indicativo) & \fonte{DS-Si4464 p.\,20--21}\\
6 & RST del DS3231M basso dopo il power-up & $t_\text{REC}$ = \SI{250}{\milli\second} tip., \SI{300}{\milli\second} max & \fonte{DS-DS3231M p.\,3}. Non influisce su I2C e RTC \fonte{DS-DS3231M p.\,4}.\\
7 & Prima scrittura possibile sulla FRAM & $t_\text{PU}$ $\geq$ \SI{250}{\micro\second} & \fonte{DS-FM25V10 p.\,19}\\
8 & Primo impulso su WDI (PC6) & entro \SI{1.12}{\second} ($t_{WD}$ min) & \fonte{LOG F-12; DS-STWD100}\\
9 & Il 5V0 arriva al VDD dello switch prima dei comandi a 3V3 & sequenza rispettata dall'hardware & \fonte{DS-HMC349A p.\,8; REV-1104 §2.2}\\
\bottomrule
\end{xltabular}

Nella tabella~\ref{tab:powerup} gli eventi sono in ordine di tempo. I valori \enquote{min/tip/max} sono minimo, tipico e massimo garantiti dal datasheet: il firmware deve sempre usare il valore peggiore (il massimo per le attese, il minimo per le scadenze). POR = reset automatico all'accensione; BOR = reset per tensione troppo bassa.

\begin{nota}[stato sicuro durante il boot]
Durante il reset e il boot dell'MCU i GPIO sono in alta impedenza. Il pull-up R702 tiene alto il gate di Q701 e la net di PC2: switch su RX, APC a zero (PA spento), LNA acceso. La radio è sotto il controllo di IC205 (modo A) \fonte{LOG rev.\ 1.41; REV-1104 §3}. Lo stato di default è quindi sicuro anche senza firmware.
\end{nota}

\chapter{Mappa completa dei pin dell'MCU}
\label{ch:pin}

\section{Come leggere le tabelle}
Le tabelle sono ricavate pin per pin dalla netlist dello schematico \fonte{SCH}. Le funzioni alternate sono quelle delle tabelle 22 e 23 del datasheet STM32L562 \fonte{DS-STM32L562 pp.\,127--129, 135}; i canali ADC vengono dalla tabella 21 \fonte{DS-STM32L562 pp.\,98--104}.

\begin{itemize}
\item \textbf{Dir}: direzione vista dall'MCU. I = ingresso, O = uscita, IO = bidirezionale, OD = open-drain, A = analogico.
\item \textbf{Pull int.}: resistenza interna consigliata (PU = pull-up, PD = pull-down, \textendash{} = nessuna).
\item \textbf{Reset}: stato del pin durante il reset e prima dell'inizializzazione. \textbf{Sicuro}: livello da impostare in caso di errore o di arresto controllato.
\item \textbf{Att.}: livello attivo della funzione.
\end{itemize}

\begin{nota}[stato dei GPIO al reset]
Dopo il reset i pin sono in alta impedenza \fonte{LOG rev.\ 1.41}. Fanno eccezione i pin di debug: PA13 (SWDIO) con pull-up, PA14 (SWCLK) con pull-down, PA15 e PB4 con pull-up; sono in funzione alternata JTAG/SWD \fonte{DS-STM32L562 p.\,126, nota 3}. Il modo esatto di reset degli altri pin (analogico) va confermato sul reference manual RM0438: \DV.
\end{nota}

% ======================================================= tabella pin usati
\begin{landscape}
\section{Pin usati}
\label{sec:pin-usati}
{\scriptsize
\rowcolors{3}{white}{sbrow}
\begin{xltabular}{\linewidth}{@{}C{6mm}L{9mm}L{30mm}L{30mm}C{6mm}C{7mm}L{24mm}C{6mm}L{34mm}Y@{}}
\caption{Pin dell'MCU U301 collegati a un componente.}\label{tab:pin}\\
\tabhead\hd{Pin}&\hd{Porta}&\hd{Net}&\hd{Funzione / periferica (AF)}&\hd{Dir}&\hd{Pull int.}&\hd{Reset $\to$ sicuro}&\hd{Att.}&\hd{Collegato a}&\hd{Note}\\
\endfirsthead
\tabhead\hd{Pin}&\hd{Porta}&\hd{Net}&\hd{Funzione / periferica (AF)}&\hd{Dir}&\hd{Pull int.}&\hd{Reset $\to$ sicuro}&\hd{Att.}&\hd{Collegato a}&\hd{Note}\\
\endhead
\multicolumn{10}{r}{\itshape continua nella pagina seguente}\\
\endfoot
\bottomrule
\endlastfoot
2 & PC13 & \net{GPIOSX\_2} & GPIO (uso da definire) & I & \textcolor{sborange}{?} & Hi-Z $\to$ ingresso & \textendash & J708 pin 13 (dock) & Pin a bassa corrente: max \SI{3}{\milli\ampere}, \SI{2}{\mega\hertz} \fonte{DS p.\,126}. \textcolor{sborange}{?} = \DV. Collegato dopo la rev.\ 1.45, vedi §\ref{sec:pc13}.\\
3 & PC14 & \net{GPIO\_PPS} & GPIO / EXTI14 (PPS dal dock) & I & \textcolor{sborange}{?} & Hi-Z $\to$ ingresso & \textcolor{sborange}{?} & J708 pin 17 (dock) & Come sopra. Polarità e livelli del PPS: \DV.\\
4 & PC15 & \net{GPIOSX\_0} & GPIO (uso da definire) & I & \textcolor{sborange}{?} & Hi-Z $\to$ ingresso & \textendash & J708 pin 14 (dock) & Come PC13.\\
5 & PH0 & \net{HSE\_IN} & RCC OSC\_IN & A & \textendash & \textendash & \textendash & Y1 \SI{24}{\mega\hertz}, C2 \SI{10}{\pico\farad} & HSE\_VALUE = 24000000 (F-06).\\
6 & PH1 & \net{HSE\_OUT} & RCC OSC\_OUT & A & \textendash & \textendash & \textendash & R1 \SI{0}{\ohm} $\to$ Y1, C1 & Quarzo CL \SI{10}{\pico\farad}, \SI{\pm10}{\ppm} \fonte{LOG 1.15-01}.\\
7 & NRST & \net{/IF/DBG.nRST} & reset & IO & \textendash & \textendash & L & R2 \SI{10}{\kilo\ohm} pull-up, C3 \SI{100}{\nano\farad}, J1.3, WDO di U403 & Il watchdog esterno resetta l'MCU da qui (F-13).\\
8 & PC0 & \net{/MCU/UHF.RF\_IRQ} & GPIO, EXTI0 fronte di discesa & I & PU & Hi-Z $\to$ ingresso & L & IC204 pin 11 nIRQ & nIRQ non richiede pull-up esterno \fonte{DS-Si4464 p.\,46}; il PU interno evita un ingresso flottante con il Si4464 in shutdown.\\
9 & PC1 & \net{/MCU/V\_TEMP\_BUCK} & ADC12\_IN2 & A & \textendash & analogico & \textendash & uscita di IC5 (TLV9051), R12 & PT1000 del buck, §\ref{sec:adc}.\\
10 & PC2 & \net{/MCU/UHF.EN\_TX\_ANT\_SW} & GPIO uscita push-pull & O & \textendash & Hi-Z (R702 $\to$ H) $\to$ \textbf{H} & L = TX & IC202 CTRL, gate Q701 (R703 \SI{0}{\ohm}, R702 \SI{100}{\kilo\ohm} a 3V3), U701 ingresso, TP702 & \textbf{Comando RF unico}. H = RX, PA spento. Sequenze: §\ref{sec:sequenze}.\\
21 & PA5 & \net{/MCU/SPI1.CLK} & SPI1\_SCK (AF5) & O & \textendash & Hi-Z $\to$ L & \textendash & IC204 SCLK & SPI modo 0, max \SI{10}{\mega\hertz} \fonte{DS-Si4464 p.\,17}.\\
22 & PA6 & \net{/MCU/SPI1.MISO} & SPI1\_MISO (AF5) & I & \textendash & Hi-Z & \textendash & IC204 SDO & \\
23 & PA7 & \net{/MCU/SPI1.MOSI} & SPI1\_MOSI (AF5) & O & \textendash & Hi-Z $\to$ L & \textendash & IC204 SDI & \\
24 & PC4 & \net{RF\_MR\_N} & GPIO uscita push-pull & O & \textendash & Hi-Z (R208 $\to$ H) $\to$ L & L & IC205 \(\overline{\mathrm{MR}}\), R208 \SI{10}{\kilo\ohm} a 3V3, R205 (DNP) verso SDN & Modo A (di serie): L = radio in reset. Modo B: logica invertita (F-05), §\ref{sec:sdn}.\\
26 & PB1 & \net{/MCU/V\_TEMP\_PA} & ADC12\_IN16 & A & \textendash & analogico & \textendash & uscita di IC1 (TLV9051), R5 & PT1000 del PA, §\ref{sec:adc}.\\
28 & PB10 & \net{/MCU/I2C2\_SCL} & I2C2\_SCL (AF4) & OD & \textendash & Hi-Z & \textendash & U404 SCL, R408 \SI{10}{\kilo\ohm} a 3V3 & RTC, indirizzo \code{0x68}.\\
29 & PB11 & \net{/MCU/I2C2\_SDA} & I2C2\_SDA (AF4) & OD & \textendash & Hi-Z & \textendash & U404 SDA, R407 \SI{10}{\kilo\ohm} a 3V3 & \\
33 & PB12 & \net{/MCU/FRAM\_CS} & GPIO uscita (CS software) & O & \textendash & Hi-Z $\to$ \textbf{H} & L & U401 /CS & \textbf{Nessun pull-up esterno}: portarlo alto per primo (F-15).\\
34 & PB13 & \net{/MCU/SPI2\_SCK} & SPI2\_SCK (AF5) & O & \textendash & Hi-Z $\to$ L & \textendash & U401 SCK & SPI modo 0 o 3, max \SI{40}{\mega\hertz} \fonte{DS-FM25V10 p.\,17}.\\
35 & PB14 & \net{/MCU/SPI2\_MISO} & SPI2\_MISO (AF5) & I & \textendash & Hi-Z & \textendash & U401 SO & \\
36 & PB15 & \net{/MCU/SPI2\_MOSI} & SPI2\_MOSI (AF5) & O & \textendash & Hi-Z $\to$ L & \textendash & U401 SI & /WP e /HOLD di U401 sono fissi a 3V3.\\
37 & PC6 & \net{/MCU/WDT\_WDI} & GPIO uscita push-pull & O & \textendash & Hi-Z $\to$ commuta & $\updownarrow$ & U403 WDI & Commutare \textbf{da software}, mai con timer o PWM (F-11). Non usare TIM3\_CH1 (AF2).\\
41 & PA8 & \net{/MCU/SPI1.CS} & GPIO uscita (NSEL software) & O & \textendash & Hi-Z $\to$ \textbf{H} & L & IC204 nSEL & \textbf{Nessun pull-up esterno}: portarlo alto per primo (F-15).\\
42 & PA9 & \net{/IF/UART1.TX} & USART1\_TX (AF7) & O & \textendash & Hi-Z $\to$ H & \textendash & U402 RXD (CH340E) & Console (F-07).\\
43 & PA10 & \net{/IF/UART1.RX} & USART1\_RX (AF7) & I & PU & Hi-Z & \textendash & U402 TXD (CH340E) & Il PU tiene la linea a riposo se il CH340E non pilota.\\
44 & PA11 & \net{/IF/CAN\_MCU1.CAN\_RX} & FDCAN1\_RX (AF9) & I & \textendash & Hi-Z & \textendash & IC501 RXD & \\
45 & PA12 & \net{/IF/CAN\_MCU1.CAN\_TX} & FDCAN1\_TX (AF9) & O & \textendash & Hi-Z $\to$ H (recessivo) & \textendash & IC501 TXD & TXD alto = recessivo.\\
46 & PA13 & \net{/IF/DBG.SWDIO} & SWDIO (AF0) & IO & PU (reset) & AF debug & \textendash & R16 \SI{22}{\ohm} $\to$ J1.2 & Non riconfigurare se serve il debug.\\
49 & PA14 & \net{/IF/DBG.SWCLK} & SWCLK (AF0) & I & PD (reset) & AF debug & \textendash & R17 \SI{22}{\ohm} $\to$ J1.4 & \\
54 & PB3 & \net{/IF/DBG.SWO} & TRACESWO (AF0) & O & \textendash & AF debug & \textendash & R11 \SI{22}{\ohm} $\to$ J1.6 & Opzionale (printf via SWO).\\
57 & PB6 & \net{/IF/I2C1.SCL} & I2C1\_SCL (AF4) & OD & \textendash & Hi-Z & \textendash & J203.3, J708.11, D202, R207 \SI{4.7}{\kilo\ohm} via JP201 & Bus esterno. JP201 aperto = nessun pull-up a bordo \fonte{LOG R-13}.\\
58 & PB7 & \net{/IF/I2C1.SDA} & I2C1\_SDA (AF4) & OD & \textendash & Hi-Z & \textendash & J203.2, J708.9, D202, R206 \SI{4.7}{\kilo\ohm} via JP201 & \\
59 & PH3 & \net{/IF/DBG.BOOT0} & BOOT0 & I & \textendash & L (R3) & H & R3 \SI{10}{\kilo\ohm} a GND, TP1 & Alto al reset = bootloader di sistema (F-07).\\
61 & PB9 & \net{/MCU/RTC\_RSTn} & GPIO ingresso & I & \textendash & Hi-Z $\to$ ingresso & L & U404 pin 4 \(\overline{\mathrm{RST}}\) & Pull-up interno \SI{50}{\kilo\ohm} nel DS3231M \fonte{DS-DS3231M p.\,7}. \textbf{Mai uscita push-pull.}\\
\end{xltabular}}
\end{landscape}

\section{Pin di alimentazione}
\begin{xltabular}{\linewidth}{@{}L{30mm}L{40mm}Y@{}}
\tabhead\hd{Pin}&\hd{Net}&\hd{Note}\\
1 (VBAT) & \net{3V3} & Nessuna batteria di backup per l'MCU.\\
13 (VDDA) & \net{3V3} & Riferimento dell'ADC (VREF+ collegato a VDDA nel package, \DV{} su RM0438/datasheet).\\
19, 32, 64 (VDD) & \net{3V3} & Disaccoppiamento C301\ldots C307 \fonte{REV-1104 §2.3}.\\
48 (VDDUSB) & \net{3V3} & USB interno dell'MCU non usato.\\
12 (VSSA), 18, 31, 47, 63 (VSS) & \net{GND} & \\
30, 62 (VDD12) & non collegati & Corretto senza SMPS esterno \fonte{REV-1104 §2.3}.\\
\bottomrule
\end{xltabular}

\section{Pin liberi}
\label{sec:pin-liberi}
Diciotto GPIO non sono collegati (simbolo di non collegato sullo schematico) \fonte{SCH}. Configurarli in modo \textbf{analogico}: è lo stato a consumo minimo e non lascia ingressi digitali flottanti.

\begin{xltabular}{\linewidth}{@{}C{10mm}L{14mm}L{50mm}Y@{}}
\tabhead\hd{Pin}&\hd{Porta}&\hd{Funzioni utili (riserva)}&\hd{Configurazione consigliata}\\
\endhead
11 & PC3 & ADC12\_IN4 & analogico\\
14 & PA0 & ADC12\_IN5, wakeup & analogico\\
15 & PA1 & ADC12\_IN6 & analogico\\
16 & PA2 & ADC12\_IN7 & analogico\\
17 & PA3 & ADC12\_IN8 & analogico\\
20 & PA4 & ADC12\_IN9 & analogico\\
25 & PB0 & ADC12\_IN15 & analogico\\
27 & PB2 & \textendash & analogico\\
38 & PC7 & \textendash & analogico\\
39 & PC8 & \textendash & analogico\\
40 & PC9 & \textendash & analogico\\
50 & PA15 & JTDI (pull-up al reset) & analogico dopo l'avvio\\
51 & PC10 & \textendash & analogico\\
52 & PC11 & \textendash & analogico\\
53 & PC12 & \textendash & analogico\\
55 & PB4 & NJTRST (pull-up al reset) & analogico dopo l'avvio\\
56 & PB5 & \textendash & analogico\\
60 & PB8 & \textendash & analogico\\
\bottomrule
\end{xltabular}
\fonte{canali ADC: DS-STM32L562 pp.\,100--103}

\section{Boot, debug e reset}
\begin{xltabular}{\linewidth}{@{}L{22mm}L{30mm}L{40mm}Y@{}}
\tabhead\hd{Segnale}&\hd{Pin MCU}&\hd{Hardware}&\hd{Uso}\\
\endhead
NRST & 7 & R2 \SI{10}{\kilo\ohm} a 3V3, C3 \SI{100}{\nano\farad}, J1 pin 3, WDO di U403 (open-drain) & Reset da debugger o da watchdog.\\
BOOT0 & 59 (PH3) & R3 \SI{10}{\kilo\ohm} a GND, TP1 & Normale: boot da flash. TP1 a 3V3 durante il reset: bootloader di sistema su USART1 (STM32CubeProgrammer) \fonte{LOG F-07}.\\
SWDIO & 46 (PA13) & R16 \SI{22}{\ohm} $\to$ J1 pin 2 & Debug SWD.\\
SWCLK & 49 (PA14) & R17 \SI{22}{\ohm} $\to$ J1 pin 4 & Debug SWD.\\
SWO & 54 (PB3) & R11 \SI{22}{\ohm} $\to$ J1 pin 6 & Trace opzionale.\\
VCC / GND & \textendash & J1 pin 1 = 3V3, pin 5 = GND & Il debugger legge il 3V3: la scheda va alimentata (EN\_DOCK).\\
\bottomrule
\end{xltabular}
Il connettore J1 è un TC2030 (Tag-Connect, solo footprint) \fonte{SCH; LOG P-15}.

\begin{attenzione}[watchdog e debug]
Con R413 non montata (configurazione di volo) il watchdog è sempre attivo. Un breakpoint, il bootloader UART o un firmware di prova che non commuta PC6 provocano un reset ogni circa \SI{1.6}{\second}. Durante lo sviluppo montare R413 (\SI{0}{\ohm}, 0402): EN alto, watchdog spento \fonte{LOG rev.\ 1.41--1.42}.
\end{attenzione}

\section{PC13, PC14 e PC15 verso il dock}
\label{sec:pc13}
\begin{verifica}[incongruenza tra schematico e log]
Nella netlist attuale (schematico salvato il 04/11/2026 alle 15:22) PC13, PC14 e PC15 sono collegati a J708 pin 13 (\net{GPIOSX\_2}), 17 (\net{GPIO\_PPS}) e 14 (\net{GPIOSX\_0}). Nella netlist della rev.\ 1.45 erano non collegati, e il log (R-12, S-09; revisione del 04/11, P15) dice che PPS e GPIOSX \enquote{non arrivano all'MCU}. La modifica non è registrata nel log. Funzione, direzione e livelli di questi segnali vanno chiariti con il team: \DV.
\end{verifica}
Fino a una decisione, configurare i tre pin come ingressi senza pull (o in modo analogico). Se in futuro diventano uscite, ricordare i limiti di questi pin: al massimo \SI{3}{\milli\ampere}, \SI{2}{\mega\hertz} con \SI{30}{\pico\farad} di carico, mai come sorgente di corrente \fonte{DS-STM32L562 p.\,126, nota 1}. PC14 e PC15 sono anche i pin dell'oscillatore LSE: sulla scheda non c'è un quarzo a \SI{32}{\kilo\hertz}, quindi l'RTC interno dell'MCU può usare solo l'LSI.

\chapter{Periferiche e bus}
\label{ch:periferiche}

\section{Clock dell'MCU}
\begin{itemize}
\item \textbf{HSE}: quarzo Y1 \SI{24}{\mega\hertz} (YXC X322524MMB4SI, CL \SI{10}{\pico\farad}, \SI{\pm10}{\ppm}) su PH0/PH1 \fonte{LOG 1.15-01}. Impostare \code{HSE\_VALUE = 24000000} \fonte{LOG F-06}. Con un valore diverso UART, CAN e timer lavorano alla frequenza sbagliata.
\item \textbf{LSE}: non presente. Per l'RTC interno usare LSI; l'ora di riferimento è il DS3231M.
\item \textbf{SYSCLK}: libera scelta del firmwarista (PLL da HSE). Le frequenze dei bus vanno scelte in modo che i prescaler SPI diano al massimo \SI{10}{\mega\hertz} su SPI1 (vedi sotto).
\end{itemize}

\begin{nota}[TrustZone]
L'STM32L562 ha TrustZone. Con TrustZone attivo i GPIO sono \enquote{secure} dopo il reset \fonte{DS-STM32L562 p.\,33}. Se non serve, lasciare l'option byte TZEN disattivato: semplifica il firmware.
\end{nota}

\section{SPI1 \textendash{} Si4464}
\begin{xltabular}{\linewidth}{@{}L{38mm}Y@{}}
\tabhead\hd{Parametro}&\hd{Valore}\\
Pin & PA5 SCK, PA6 MISO, PA7 MOSI (AF5); PA8 NSEL come GPIO.\\
Modo & CPOL = 0, CPHA = 0 (modo 0). Il Si4464 campiona SDI sul fronte di salita e cambia SDO sul fronte di discesa; l'host legge SDO sul fronte di salita \fonte{DS-Si4464 p.\,17}.\\
Ordine dei bit & MSB per primo, parole da 8 bit.\\
Clock massimo & \SI{10}{\mega\hertz} \fonte{DS-Si4464 p.\,17}.\\
Tempi minimi & $t_\text{SS}$ \SI{20}{\nano\second}, $t_\text{SH}$ \SI{50}{\nano\second}, $t_\text{SW}$ (nSEL alto tra due comandi) \SI{80}{\nano\second} \fonte{DS-Si4464 p.\,17, tab.\,8}.\\
Chip select & Software (GPIO). Una transazione = nSEL basso, comando, argomenti, nSEL alto.\\
Protocollo & Dopo ogni comando leggere CTS con \code{READ\_CMD\_BUFF} finché vale \code{0xFF}; tempo tipico \SI{20}{\micro\second} \fonte{DS-Si4464 p.\,17}. I registri \enquote{Fast Response} (comandi \code{0x50}, \code{0x51}, \code{0x53}, \code{0x57}) si leggono senza CTS \fonte{DS-Si4464 p.\,19}.\\
\bottomrule
\end{xltabular}

\section{SPI2 \textendash{} FRAM FM25V10}
\begin{xltabular}{\linewidth}{@{}L{38mm}Y@{}}
\tabhead\hd{Parametro}&\hd{Valore}\\
Pin & PB13 SCK, PB14 MISO, PB15 MOSI (AF5); PB12 /CS come GPIO.\\
Modo & Modo 0 o modo 3 (il chip riconosce il modo dal livello di SCK quando /CS scende) \fonte{DS-FM25V10 pp.\,5--6}.\\
Clock massimo & \SI{40}{\mega\hertz} con VDD \SIrange{2.7}{3.6}{\volt} \fonte{DS-FM25V10 p.\,17}.\\
Organizzazione & \num{131072} $\times$ 8 bit, indirizzo di 3 byte (17 bit utili, i 7 bit alti \enquote{don't care}) \fonte{DS-FM25V10 pp.\,4--5}.\\
Scrittura & Nessun ritardo di scrittura (\enquote{NoDelay}). Prima di ogni WRITE o WRSR serve WREN \fonte{DS-FM25V10 p.\,6}.\\
/WP, /HOLD & Fissi a 3V3: la protezione hardware dello status register non è usabile.\\
Avvio & Nessun accesso prima di $t_\text{PU}$ = \SI{250}{\micro\second} dal power-up \fonte{DS-FM25V10 p.\,19}.\\
\bottomrule
\end{xltabular}

\begin{xltabular}{\linewidth}{@{}L{18mm}L{16mm}Y@{}}
\caption{Codici operativi della FM25V10 \fonte{DS-FM25V10 p.\,6}.}\\
\tabhead\hd{Comando}&\hd{Codice}&\hd{Funzione}\\
WREN & \code{0x06} & abilita la scrittura (latch WEL)\\
WRDI & \code{0x04} & disabilita la scrittura\\
RDSR & \code{0x05} & legge lo status register\\
WRSR & \code{0x01} & scrive lo status register (BP1, BP0, WPEN)\\
READ & \code{0x03} & lettura\\
FSTRD & \code{0x0B} & lettura veloce (con byte dummy)\\
WRITE & \code{0x02} & scrittura\\
SLEEP & \code{0xB9} & sleep; risveglio con /CS, $t_\text{REC}$ \SI{400}{\micro\second} max \fonte{DS-FM25V10 pp.\,11, 19}\\
RDID & \code{0x9F} & Device ID, 9 byte: \code{7F 7F 7F 7F 7F 7F C2 24 00} \fonte{DS-FM25V10 p.\,11}\\
\bottomrule
\end{xltabular}

\begin{regola}[F-15: chip select all'avvio]
All'avvio configurare il prima possibile PB12 (CS della FRAM) e PA8 (NSEL del Si4464) come uscite a livello \textbf{alto}, prima di qualsiasi altra inizializzazione SPI \fonte{LOG F-15}. Le due linee non hanno pull-up esterni (decisione P-12, \enquote{Lascia com'è}).
\end{regola}

\section{I\textsuperscript{2}C}
\begin{xltabular}{\linewidth}{@{}L{20mm}L{28mm}L{40mm}Y@{}}
\tabhead\hd{Bus}&\hd{Pin}&\hd{Pull-up}&\hd{Dispositivi}\\
I2C2 & PB10 SCL, PB11 SDA (AF4) & R408, R407 \SI{10}{\kilo\ohm} a 3V3, sempre montati & U404 DS3231M, indirizzo 7 bit \code{0x68} (byte \code{0xD0}/\code{0xD1}) \fonte{DS-DS3231M p.\,16}\\
I2C1 & PB6 SCL, PB7 SDA (AF4) & R207, R206 \SI{4.7}{\kilo\ohm} a 3V3 \textbf{solo con JP201 chiuso} & bus esterno verso J203 e dock J708 (pin 9, 11). Protezione ESD D202 (PESD5V0L2BT)\\
\bottomrule
\end{xltabular}

\begin{itemize}
\item \textbf{Velocità I2C2.} Il DS3231M supporta \SI{400}{\kilo\hertz} \fonte{DS-DS3231M p.\,1}. Con pull-up da \SI{10}{\kilo\ohm} il tempo di salita dipende dalla capacità del bus: \SI{100}{\kilo\hertz} è la scelta prudente; \SI{400}{\kilo\hertz} va confermato misurando il tempo di salita (limite \SI{300}{\nano\second} in fast mode) \DV.
\item \textbf{JP201 (I2C1).} Aperto di serie: i pull-up li fornisce il dock. Si chiude solo sulla scheda che deve dare i pull-up al bus \fonte{LOG R-13}. I pull-up interni dell'STM32 non sostituiscono quelli esterni. Se il bus resta a livello basso senza dock, il firmware deve gestire il timeout e non bloccarsi.
\end{itemize}

\section{UART e USB}
\begin{itemize}
\item \textbf{USART1} su PA9 (TX) e PA10 (RX), AF7, collegata al CH340E U402 \fonte{LOG F-07}. Il CH340E supporta da \SI{50}{bps} a \SI{2}{Mbps} \fonte{DS-CH340E p.\,1}. RTS, CTS e TNOW del CH340E non sono collegati: niente controllo di flusso.
\item Il CH340E è alimentato dal 3V3 (VCC e V3 uniti) \fonte{REV-1104 §2.4}: con la scheda spenta la porta USB non compare sul PC.
\item L'USB interno dell'STM32 non è usato (PA11/PA12 vanno al CAN).
\item \textbf{Bootloader}: TP1 (BOOT0) alto al reset, poi STM32CubeProgrammer su UART \fonte{LOG F-07}. Con il watchdog attivo serve R413 montata (capitolo~\ref{ch:supervisione}).
\end{itemize}

\section{CAN}
\begin{itemize}
\item \textbf{FDCAN1} su PA11 (RX) e PA12 (TX), AF9 \fonte{DS-STM32L562 p.\,135}. Transceiver TCAN3414 (IC501) in modo normale fisso: STB e SHDN sono a GND \fonte{LOG R-17; REV-1104 §2.4}.
\item CAN classico e CAN FD fino a \SI{8}{Mbps} \fonte{DS-TCAN3414 p.\,1}.
\item Timeout di dominante su TXD: \SIrange{1.2}{4.0}{\milli\second} \fonte{DS-TCAN3414 p.\,9}. Un TXD basso più a lungo disabilita il driver: limita il bitrate minimo utilizzabile.
\item Terminazione \SI{120}{\ohm} (R27) inserita solo con JP6 chiuso \fonte{SCH}. Linee verso J202 e J708 (pin 10 CANL, 12 CANH), protezione ESD D201.
\end{itemize}

\section{ADC: sensori di temperatura PT1000}
\label{sec:adc}
\begin{xltabular}{\linewidth}{@{}L{30mm}L{22mm}L{24mm}Y@{}}
\tabhead\hd{Misura}&\hd{Pin}&\hd{Canale}&\hd{Catena}\\
Temperatura PA & PB1 (pin 26) & ADC12\_IN16 & R\_RTD2 PT1000, R4, IC1 TLV9051, R5/R6/R7\\
Temperatura buck & PC1 (pin 9) & ADC12\_IN2 & R\_RTD1 PT1000, R8, IC5 TLV9051, R12/R9/R10\\
\bottomrule
\end{xltabular}
\fonte{LOG F-08, S-06; DS-STM32L562 pp.\,98, 104}

\subsection{Circuito}
Ogni PT1000 forma un partitore con una resistenza da \SI{3.16}{\kilo\ohm} verso 3V3. L'operazionale amplifica la differenza rispetto a un riferimento ricavato dallo stesso 3V3. Le resistenze sono allo \SI{0.1}{\percent} \fonte{LOG rev.\ 1.18}.

\begin{center}
\begin{circuitikz}[scale=0.85,transform shape,font=\small]
\draw (0,4) node[above]{3V3} to[R,l=$R_4$ \SI{3.16}{\kilo\ohm},o-*] (0,2) coordinate(vp)
      to[R,l=$R_t$ (PT1000)] (0,0) node[ground]{};
\draw (vp) -- (2,2) node[op amp,anchor=+,yscale=-1] (oa) {};
\draw (oa.-) -- ++(-0.6,0) coordinate(vm);
\draw (vm) to[R,l_=$R_7$ \SI{19.6}{\kilo\ohm},*-] ++(0,-2.2) node[ground]{};
\draw (vm) |- ++(0,1.6) coordinate(top);
\draw (top) to[R,l=$R_6$ \SI{63.4}{\kilo\ohm}] ++(-2.8,0) node[left]{3V3};
\draw (top) to[R,l=$R_5$ \SI{105}{\kilo\ohm}] (top -| oa.out) -- (oa.out);
\draw (oa.out) to[short,*-o] ++(1.2,0) node[right]{ADC (PB1 / PC1)};
\node[font=\scriptsize] at ($(oa)+(0.2,-1.1)$) {TLV9051};
\end{circuitikz}
\end{center}
Valori per la catena del PA (IC1); la catena del buck (IC5) usa R8, R12, R9, R10 con gli stessi valori \fonte{SCH}.

\subsection{Formule di conversione}
Con $G = 1 + R_5/R_6 + R_5/R_7$ e $k = R_5/R_6$:
\begin{align}
\frac{V_\text{out}}{V_\text{DDA}} &= G\,\frac{R_t}{R_4+R_t} - k,
\qquad G = 1+\frac{105}{63{,}4}+\frac{105}{19{,}6} = 8{,}0133,\quad k = \frac{105}{63{,}4}=1{,}65615\\
\intertext{L'ADC usa VDDA (= 3V3) come riferimento, quindi la misura è raziometrica: con un codice $N$ a 12 bit, $r = N/4095$,}
x &= \frac{r + k}{G},\qquad R_t = R_4\,\frac{x}{1-x}\\
\intertext{e la temperatura si ricava dalla PT1000 (IEC 60751, $R_0 = \SI{1000}{\ohm}$, $A = \num{3.9083e-3}$, $B = \num{-5.775e-7}$):}
T &= \frac{-A + \sqrt{A^2 - 4B\,(1 - R_t/R_0)}}{2B}.
\end{align}
Il termine $C$ della norma, valido sotto \SI{0}{\celsius}, sposta il risultato di circa \SI{0.01}{\celsius} a \SI{-40}{\celsius} (calcolo: $R_0 C (T-100) T^3 \approx \SI{0.04}{\ohm}$ contro $\approx \SI{3.9}{\ohm\per\celsius}$) e si può trascurare.

\textbf{Verifiche del calcolo}: $V_\text{out}=0$ per $R_t = \SI{823}{\ohm}$ ($\approx\SI{-46}{\celsius}$) e $V_\text{out} = V_\text{DDA}$ per $R_t = \SI{1567}{\ohm}$ ($\approx\SI{+147}{\celsius}$), in accordo con la revisione \fonte{REV-1104 §2.5}. A \SI{-40}{\celsius} l'uscita vale \SI{0.10}{\volt}, a \SI{+125}{\celsius} \SI{2.97}{\volt} \fonte{LOG R-18}.

\begin{center}
\begin{tikzpicture}
\begin{axis}[width=0.78\linewidth,height=6cm,grid=both,grid style={gray!20},
  xlabel={Temperatura [\si{\celsius}]},ylabel={$V_\text{out}$ [\si{\volt}]},
  xmin=-50,xmax=150,ymin=0,ymax=3.3,
  ytick={0,0.5,1,1.5,2,2.5,3,3.3},
  yticklabel style={/pgf/number format/use comma},
  tick label style={font=\footnotesize},label style={font=\small\sffamily},
  legend style={font=\footnotesize,at={(0.03,0.97)},anchor=north west}]
\addplot[very thick,sbblue,domain=-46:147,samples=80]
  {3.3*(8.0133*(1000*(1+3.9083e-3*x-5.775e-7*x^2))/(3160+1000*(1+3.9083e-3*x-5.775e-7*x^2))-1.65615)};
\addlegendentry{$V_\text{out}$ con 3V3 = \SI{3.30}{\volt}}
\addplot[only marks,mark=*,sbred] coordinates {(-40,0.102) (125,2.967)};
\addlegendentry{punti del log R-18}
\end{axis}
\end{tikzpicture}
\end{center}

La sensibilità non è costante: \SI{20.6}{\milli\volt\per\celsius} a \SI{-40}{\celsius}, \SI{18.3}{} a \SI{15}{\celsius}, \SI{14.6}{\milli\volt\per\celsius} a \SI{125}{\celsius} (calcolo dalle formule sopra). La media fra \SI{-40}{} e \SI{125}{\celsius} è \SI{17.4}{\milli\volt\per\celsius}, in accordo con il valore $\approx$\SI{17.3}{\milli\volt\per\celsius} del log \fonte{LOG R-18}. Un LSB a 12 bit (\SI{0.81}{\milli\volt}) vale quindi circa \SIrange{0.04}{0.06}{\celsius}. Corrente nella PT1000: $3{,}3/4{,}16\,\text{k} \approx \SI{0.79}{\milli\ampere}$ \fonte{REV-1104 §2.5}.

\begin{regola}[F-08: calibrazione]
Misura raziometrica su VDDA. Fare una calibrazione a un punto per canale e salvare l'offset (per esempio in FRAM). Senza calibrazione l'offset del TLV9051 vale circa \SI{\pm0.9}{\celsius} \fonte{LOG F-08, R-18}.
\end{regola}

\begin{nota}[impostazioni ADC]
Le uscite dei TLV9051 hanno bassa impedenza: non servono tempi di campionamento lunghi \fonte{LOG §29}. L'oversampling hardware dell'ADC riduce il rumore. Il tempo di campionamento minimo per l'impedenza di sorgente è nella tabella 110 del datasheet \fonte{DS-STM32L562 p.\,246}.
\end{nota}

\section{Timer}
Nessuna funzione hardware richiede un timer. Il firmware può usare liberamente SysTick e i timer per la sua temporizzazione, con due vincoli:
\begin{itemize}
\item \textbf{PC6 (WDI) non va mai pilotato da un timer} (TIM3\_CH1 o PWM): il refresh del watchdog deve dipendere dal ciclo principale \fonte{LOG F-11}.
\item Le attese delle sequenze RF (\SI{10}{\micro\second}, capitolo~\ref{ch:rf}) richiedono una base dei tempi con risoluzione di almeno \SI{1}{\micro\second} (per esempio il contatore DWT\_CYCCNT del Cortex-M33 o un timer a \SI{1}{\mega\hertz}).
\end{itemize}

\section{Interrupt}
\begin{xltabular}{\linewidth}{@{}L{30mm}L{26mm}L{26mm}Y@{}}
\tabhead\hd{Sorgente}&\hd{Pin}&\hd{Linea}&\hd{Configurazione}\\
Si4464 nIRQ & PC0 & EXTI0 & Fronte di discesa, pull-up interno. Nella ISR leggere e cancellare lo stato con \code{GET\_INT\_STATUS} (\code{0x20}) \DV{} sulla API della revisione B1.\\
PPS dal dock & PC14 & EXTI14 & Solo se il segnale verrà usato (§\ref{sec:pc13}), \DV.\\
FDCAN1, USART1, I2C, SPI & \textendash & interrupt di periferica & Liberi.\\
\bottomrule
\end{xltabular}

Segnali che \textbf{non} arrivano all'MCU: INT/SQW del DS3231M (solo pull-up R411) e uscita 32\,kHz del DS3231M \fonte{LOG S-09, R-12; REV-1104 P15}; GPIO0\ldots GPIO3 e TXRAMP del Si4464 \fonte{SCH}; ST del TPS2121 e PG del buck.

\chapter{Parte RF}
\label{ch:rf}

Questo è il capitolo più importante. La catena RF non ha protezioni hardware contro alcuni errori di sequenza: \textbf{la protezione del commutatore e del PA è affidata al firmware} \fonte{LOG P-02, R-04}.

\section{Architettura TX e RX}
\label{sec:rf-arch}

\subsection{Catena di trasmissione}
Si4464 (IC204, uscita TX in classe E, match AN648) $\to$ attenuatore a $\pi$ R245/R246/R247+R248 ($\approx$\SI{10.8}{dB}) $\to$ PA QPA9510 (IC2) $\to$ rete d'uscita del PA $\to$ filtro passa-basso FL201 (LFCG-490+) $\to$ porta RF2 dello switch HMC349A (IC202) $\to$ porta comune RFC $\to$ MMCX J201 \fonte{SCH; LOG §41; REV-1104 §2.2}.

\subsection{Catena di ricezione}
MMCX $\to$ RFC $\to$ porta RF1 dello switch $\to$ filtro passa-basso C209--L201--C210 $\to$ LNA TQP3M9036 (Q201) $\to$ match RX (AN643, 4 elementi) $\to$ ingressi differenziali RXp/RXn del Si4464 \fonte{SCH; LOG rev.\ 1.45, P11}.

\begin{nota}[nota sullo schematico]
Vicino al filtro RX lo schematico dice \enquote{435MHz filtro passa banda}, ma il filtro è un passa-basso. Il satellite è solo UHF: la nota va ignorata \fonte{LOG rev.\ 1.45, P11}.
\end{nota}

\subsection{Il comando unico PC2}
PC2 (net \net{/MCU/UHF.EN\_TX\_ANT\_SW}) pilota contemporaneamente tre cose \fonte{LOG §41; SCH}:
\begin{enumerate}
\item il pin CTRL dello switch HMC349A (EN dello switch è a GND);
\item il gate del P-MOS Q701 (FDN308P) tramite R703 \SI{0}{\ohm}. Il source di Q701 è sul 3V3, il pull-up R702 \SI{100}{\kilo\ohm} tiene il gate alto. Il drain alimenta APC1/APC2 del PA tramite il partitore R242 \SI{33}{\ohm} / R243 \SI{220}{\ohm}, con C252 \SI{10}{\nano\farad};
\item l'ingresso dell'invertitore U701 (74AUP2G14), la cui uscita \net{EN\_RX\_LNA} va al pin Shut Down dell'LNA tramite R202 \SI{0}{\ohm}.
\end{enumerate}

\begin{figure}[htbp]
\centering
\begin{circuitikz}[scale=0.82,transform shape,font=\small]
 % PC2
 \draw (0,0) node[left]{PC2 (U301 pin 10)} to[short,o-*] (1.5,0) coordinate(n);
 \draw (n) -- (1.5,-0.0) node[below right,font=\scriptsize]{TP702};
 % verso switch
 \draw (n) |- (3,2.5) node[right]{CTRL HMC349A (IC202): H = RF1/RX, L = RF2/TX};
 % verso Q701
 \draw (n) to[R,l=R703 \SI{0}{\ohm},-*] (4.5,0) coordinate(g);
 \draw (g) to[R,l_=R702 \SI{100}{\kilo\ohm}] (4.5,2) node[above]{3V3};
 \draw (g) -- (5.4,0) node[pmos,anchor=G,xscale=1] (q) {};
 \node[right,font=\scriptsize] at ($(q.center)+(0.3,0.0)$) {Q701};
 \draw (q.S) -- ++(0,0.5) node[above]{3V3};
 \draw (q.D) to[R,l=R242 \SI{33}{\ohm},-*] ++(0,-2) coordinate(apc);
 \draw (apc) to[R,l=R243 \SI{220}{\ohm}] ++(0,-2) node[ground]{};
 \draw (apc) -- ++(1.8,0) coordinate(c1) to[C,l=C252 \SI{10}{\nano\farad}] ++(0,-2) node[ground]{};
 \draw (c1) -- ++(1.4,0) node[right]{APC1/APC2 del PA (IC2)};
 % verso U701
 \draw (n) |- (3,-3.2) node[not port,anchor=in 1,scale=0.9] (inv) {};
 \node[below,font=\scriptsize] at (inv.south) {U701 (74AUP2G14)};
 \draw (inv.out) to[R,l=R202 \SI{0}{\ohm}] ++(2.2,0) node[right]{Shut Down LNA (Q201): H = spento};
\end{circuitikz}
\caption{Il comando PC2 e i tre carichi. Schema semplificato dalla netlist.}
\label{fig:pc2}
\end{figure}

\textbf{Conseguenze.} PC2 \textbf{basso} = switch su TX, APC alimentata (PA acceso), LNA spento. PC2 \textbf{alto} (o in alta impedenza, grazie a R702) = switch su RX, APC a zero, LNA acceso \fonte{REV-1104 §2.2}.

I livelli logici sono compatibili: HMC349A $V_\text{INH}$ \SI{2}{\volt} min, TQP3M9036 spento da \SI{2.5}{\volt} a VDD e acceso da \SIrange{0}{0.4}{\volt} \fonte{REV-1104 §2.2; DS-TQP3M9036 p.\,2}. Il percorso dello switch segue la tabella di verità del datasheet: con EN = L, CTRL = H collega RF1 (RX), CTRL = L collega RF2 (TX) \fonte{DS-HMC349A p.\,8, tab.\,4}.

\section{Tabelle di verità}
\label{sec:verita}

\subsection{Stati della catena RF}
{\small
\rowcolors{2}{white}{sbrow}
\begin{xltabular}{\linewidth}{@{}L{21mm}C{13mm}L{15mm}L{11mm}L{16mm}L{15mm}L{13mm}Y@{}}
\caption{Stati della catena RF in funzione di PC2 e dello stato del Si4464.}\label{tab:verita}\\
\tabhead\hd{Stato}&\hd{PC2}&\hd{Switch}&\hd{Q701}&\hd{APC}&\hd{PA}&\hd{LNA}&\hd{Si4464}\\
\endfirsthead
\tabhead\hd{Stato}&\hd{PC2}&\hd{Switch}&\hd{Q701}&\hd{APC}&\hd{PA}&\hd{LNA}&\hd{Si4464}\\
\endhead
Reset / boot MCU & Hi-Z $\to$ H (R702) & RF1 (RX) & OFF & \SI{0}{\volt} & standby & acceso & controllato da IC205 (modo A)\\
OFF (radio spenta) & H & RF1 (RX) & OFF & \SI{0}{\volt} & standby & acceso & shutdown (SDN = H)\\
IDLE / READY & H & RF1 (RX) & OFF & \SI{0}{\volt} & standby & acceso & READY o SLEEP\\
RX & H & RF1 (RX) & OFF & \SI{0}{\volt} & standby & acceso & RX\\
TX armato & L & RF2 (TX) & ON & \SIrange{2.7}{2.87}{\volt} & acceso, senza pilotaggio & spento & READY (\textbf{non} in TX)\\
TX & L & RF2 (TX) & ON & \SIrange{2.7}{2.87}{\volt} & acceso & spento & TX\\
\rowcolor{sbredbg}\textbf{VIETATO} & commuta & \textendash & \textendash & \textendash & \textendash & \textendash & in TX o appena uscito da TX\\
\bottomrule
\end{xltabular}}
\fonte{SCH; LOG §41, F-02, F-14; REV-1104 §2.2; APC: LOG rev.\ 1.45}

\begin{nota}[consumi da ricordare]
\begin{itemize}
\item L'LNA non si può spegnere da solo: con PC2 alto è sempre acceso, anche con la radio in shutdown. Assorbe circa \SI{45}{\milli\ampere} a \SI{3.3}{\volt} \fonte{REV-1104 §2.2, DS-TQP3M9036 p.\,7}.
\item Con PC2 basso il PA, anche senza segnale, assorbe la corrente di riposo: \SI{208}{\milli\ampere} tipici a \SI{3.6}{\volt} \fonte{DS-QPA9510 p.\,3}. Tenere PC2 basso solo per la durata della trasmissione.
\end{itemize}
\end{nota}

\subsection{Shutdown e reset del Si4464 (SDN, \texorpdfstring{$\overline{\text{MR}}$}{MR})}
\label{sec:sdn}
Il pin SDN del Si4464 è alto = radio spenta \fonte{DS-Si4464 p.\,46}. È pilotato dal supervisore IC205 (TPS3825-33, uscita RESET push-pull attiva alta) tramite R204 \SI{0}{\ohm}. PC4 (net \net{RF\_MR\_N}) va all'ingresso \(\overline{\mathrm{MR}}\) di IC205, con pull-up R208 \SI{10}{\kilo\ohm}. In laboratorio si può escludere IC205 (modo B): si toglie R204 e si monta R205 \SI{0}{\ohm}, che collega PC4 direttamente a SDN \fonte{LOG §46, R-19}.

{\small
\rowcolors{2}{white}{sbrow}
\begin{xltabular}{\linewidth}{@{}L{30mm}L{30mm}L{32mm}Y@{}}
\caption{Comando della radio con PC4 nei due modi di montaggio \fonte{LOG §46}.}\label{tab:sdn}\\
\tabhead\hd{Condizione}&\hd{PC4}&\hd{SDN Si4464}&\hd{Effetto}\\
\endfirsthead
\tabhead\hd{Condizione}&\hd{PC4}&\hd{SDN Si4464}&\hd{Effetto}\\
\endhead
\multicolumn{4}{@{}l}{\textbf{Modo A (di serie): R204 montata, R205 DNP}}\\
3V3 sotto soglia & qualsiasi & H & radio spenta (soglia \SI{2.88}{}/\SI{2.93}{}/\SI{3.00}{\volt})\\
reset radio & L (almeno \SI{1}{\micro\second}) & H, poi per $t_d$ dopo il rilascio & radio spenta, registri persi\\
normale & H o Hi-Z (R208) & L dopo $t_d$ = \SIrange{120}{300}{\milli\second} & radio attiva; D203 acceso\\
\multicolumn{4}{@{}l}{\textbf{Modo B (ripiego): R204 rimossa, R205 = \SI{0}{\ohm}}}\\
reset / boot MCU & Hi-Z (R208 $\to$ H) & H & radio spenta\\
radio spenta & H & H & radio spenta\\
radio attiva & L & L & radio attiva; D203 non indica più lo stato della radio\\
\bottomrule
\end{xltabular}}
\fonte{tempi: DS-TPS3825 pp.\,6--8; DS-Si4464 p.\,21}

\begin{regola}[F-05: reset del Si4464]
\begin{itemize}
\item Il firmware deve sapere com'è montata la scheda: costante di compilazione \code{MODO\_A} / \code{MODO\_B} \fonte{LOG F-05}.
\item \textbf{Modo A}: PC4 basso = reset. Dopo il rilascio di PC4 (o all'accensione) aspettare il ritardo del supervisore, poi il POR del Si4464, poi controllare CTS prima di mandare POWER\_UP.
\item \textbf{Modo B}: tenere SDN (= PC4) alto almeno \SI{10}{\micro\second} \fonte{DS-Si4464 p.\,21}, poi PC4 basso, attendere il POR, controllare CTS.
\end{itemize}
\end{regola}

\begin{verifica}[ritardo del supervisore]
La nota F-05 e la nota sullo schematico indicano \SI{200}{\milli\second}. È il valore tipico: il datasheet dà $t_d$ = \SI{120}{\milli\second} min, \SI{200}{\milli\second} tip., \SI{300}{\milli\second} max \fonte{DS-TPS3825 p.\,8}. Il firmware deve aspettare almeno \SI{300}{\milli\second}, oppure fare polling di CTS con un timeout di almeno \SI{300}{\milli\second} più il POR. Il tempo esatto del POR del Si4464 dopo SDN basso non è stato trovato nel datasheet (riferimento indicativo: shutdown$\to$TX/RX \SI{15}{\milli\second} \fonte{DS-Si4464 p.\,20}): \DV{} su AN625/API.
\end{verifica}

\section{Sequenze di commutazione}
\label{sec:sequenze}

\subsection{Tempi dei componenti}
\begin{xltabular}{\linewidth}{@{}L{55mm}L{34mm}Y@{}}
\tabhead\hd{Grandezza}&\hd{Valore}&\hd{Fonte}\\
Commutazione dello switch HMC349A ($t_\text{ON}$, $t_\text{OFF}$) & \SI{160}{\nano\second} & \fonte{DS-HMC349A p.\,3}\\
Accensione/spegnimento del PA da APC & \SI{175}{\nano\second} \DV & \fonte{DS-QPA9510 p.\,3}; la revisione del 04/11 riporta \SI{1.7}{\micro\second} \fonte{REV-1104 P04}\\
Salita dell'APC (Q701 ON) & $\tau = C_{252}\,(R_{242}\parallel R_{243}) = \SI{10}{\nano\farad}\times\SI{28.7}{\ohm} = \SI{0.29}{\micro\second}$ & calcolo\\
Discesa dell'APC (Q701 OFF) & $\tau \leq C_{252}R_{243} = \SI{2.2}{\micro\second}$; sotto \SI{0.5}{\volt} (\enquote{OFF}) dopo $\tau\ln(2{,}87/0{,}5) = \SI{3.8}{\micro\second}$ & calcolo; soglia OFF \fonte{DS-QPA9510 p.\,3}\\
LNA: accensione / spegnimento & \SI{1}{\micro\second} / \SI{0.5}{\micro\second} & \fonte{DS-TQP3M9036 p.\,2}\\
Invertitore U701 & $t_\text{pd}\leq$ \SI{4.3}{\nano\second} & \fonte{DS-SN74AUP2G14 p.\,1}\\
\bottomrule
\end{xltabular}
L'APC si scarica più lentamente di quanto lo switch commuti: se PC2 sale mentre il PA è ancora pilotato, per alcuni microsecondi il PA spinge circa \SI{+31}{\dBm} su una porta ormai terminata, oltre il massimo assoluto di \SI{26.5}{\dBm} \fonte{LOG P-02}.

\subsection{Sequenza RX\texorpdfstring{$\to$}{->}TX}
\begin{enumerate}
\item Verificare che il Si4464 \textbf{non} sia in TX (stato READY o RX). Se è in RX, portarlo in READY con \code{CHANGE\_STATE}.
\item Verificare che PA\_PWR\_LVL sia al valore F-01 (paragrafo~\ref{sec:limiti}).
\item Portare PC2 \textbf{basso}.
\item Attendere \textbf{almeno \SI{10}{\micro\second}}: switch, APC e LNA si assestano \fonte{REV-1104 P04; LOG F-02}.
\item Inviare \code{START\_TX}. Il Si4464 fa la sua rampa di potenza interna.
\item Attendere la fine del pacchetto (interrupt PACKET\_SENT su nIRQ). Impostare lo stato finale del TX a READY.
\end{enumerate}

\subsection{Sequenza TX\texorpdfstring{$\to$}{->}RX}
\begin{enumerate}
\item Attendere la fine del pacchetto (o forzare l'uscita da TX con \code{CHANGE\_STATE} a READY).
\item Controllare che il Si4464 sia davvero fuori da TX (stato del chip o registro Fast Response).
\item Attendere \textbf{almeno \SI{10}{\micro\second}} \fonte{LOG F-14}.
\item Solo ora portare PC2 \textbf{alto}.
\item Attendere che l'LNA si accenda (\SI{1}{\micro\second}) e l'APC si scarichi (\SI{3.8}{\micro\second}); per uniformità usare ancora \SI{10}{\micro\second}.
\item \code{START\_RX}.
\item Controllare che il Si4464 risponda ancora (F-10, paragrafo~\ref{sec:rf-supervisore}).
\end{enumerate}

\begin{attenzione}[F-14: mai commutare PC2 durante la trasmissione]
Prima si ferma la trasmissione (Si4464 fuori da TX), si attendono almeno \SI{10}{\micro\second}, solo dopo si commuta PC2 verso RX. \textbf{Mai commutare PC2 durante la trasmissione o subito dopo} \fonte{LOG F-14}. Senza questa regola il commutatore può superare il massimo assoluto di \SI{26.5}{\dBm} sulla porta terminata a ogni commutazione \fonte{LOG P-02; DS-HMC349A p.\,4}. Il limite in hot switching è \SI{30}{\dBm}.
\end{attenzione}

\begin{verifica}[testo della nota F-14]
La nota F-14 chiede di avere \enquote{PA spento/APC a zero} prima di commutare PC2. Con l'hardware attuale l'APC è comandata proprio da PC2: non si può azzerare l'APC prima di commutare PC2. L'unica cosa che il firmware può fare prima è togliere il pilotaggio al PA, cioè portare il Si4464 fuori da TX. Questo manuale interpreta F-14 in questo modo; l'interpretazione va confermata dal team.
\end{verifica}

\subsection{Diagrammi temporali}
\begin{figure}[htbp]
\centering
\begin{tikztimingtable}[timing/slope=0.12,timing/coldist=2pt,xscale=1.3,yscale=1.0,
   timing/d/background/.style={fill=sblight},font=\sffamily\scriptsize]
 Stato Si4464 & 3D{READY} 6D{READY} 2D{TX tune} 7D{TX (pacchetto)} 2D{READY}\\
 PC2          & 3H 17L \\
 Switch       & 3.1D{RF1 (RX)} 16.9D{RF2 (TX)}\\
 APC          & 3L 0.5U 16.5H\\
 LNA          & 3.1H 0.4U 16.5L\\
 RF all'MMCX  & 11L 1U 6H 1U 1L\\
 nIRQ (PC0)   & 19H 0.5L 0.5H\\
\extracode
 \tablerules
 \begin{pgfonlayer}{background}
  \vertlines[sbgray!60,densely dashed]{3,9,19}
 \end{pgfonlayer}
 \node[anchor=south,font=\sffamily\scriptsize] at (3,0.6) {PC2 $\downarrow$};
 \node[anchor=south,font=\sffamily\scriptsize] at (9,0.6) {START\_TX};
 \node[anchor=south,font=\sffamily\scriptsize] at (19,0.6) {PACKET\_SENT};
 \draw[{Stealth}-{Stealth},sbred,thick] (3,-10.8) -- node[below,font=\sffamily\scriptsize]{$\geq\SI{10}{\micro\second}$} (9,-10.8);
\end{tikztimingtable}
\caption{Sequenza RX$\to$TX. Non in scala. U = transitorio (switch \SI{160}{\nano\second}, APC $\tau$ \SI{0.29}{\micro\second}, LNA OFF \SI{0.5}{\micro\second}).}
\label{fig:rx2tx}
\end{figure}

\begin{figure}[htbp]
\centering
\begin{tikztimingtable}[timing/slope=0.12,timing/coldist=2pt,xscale=1.3,yscale=1.0,
   timing/d/background/.style={fill=sblight},font=\sffamily\scriptsize]
 Stato Si4464 & 4D{TX (fine pacchetto)} 7D{READY} 3D{READY} 6D{RX}\\
 RF all'MMCX  & 3H 1U 16L\\
 nIRQ (PC0)   & 4H 0.5L 15.5H\\
 PC2          & 11L 9H\\
 Switch       & 11.1D{RF2 (TX)} 8.9D{RF1 (RX)}\\
 APC          & 11H 2U 7L\\
 LNA          & 11L 0.8U 8.2H\\
\extracode
 \tablerules
 \begin{pgfonlayer}{background}
  \vertlines[sbgray!60,densely dashed]{4,11,14}
 \end{pgfonlayer}
 \node[anchor=south,font=\sffamily\scriptsize] at (4,0.6) {PACKET\_SENT};
 \node[anchor=south,font=\sffamily\scriptsize] at (11,0.6) {PC2 $\uparrow$};
 \node[anchor=south,font=\sffamily\scriptsize] at (14,0.6) {START\_RX};
 \draw[{Stealth}-{Stealth},sbred,thick] (4,-10.8) -- node[below,font=\sffamily\scriptsize]{$\geq\SI{10}{\micro\second}$ (F-14)} (11,-10.8);
 \draw[{Stealth}-{Stealth},sbgreen,thick] (11,-10.8) -- node[below,font=\sffamily\scriptsize]{$\geq\SI{10}{\micro\second}$} (14,-10.8);
\end{tikztimingtable}
\caption{Sequenza TX$\to$RX. Non in scala. La discesa dell'APC ($\tau\leq\SI{2.2}{\micro\second}$) avviene quando il PA non è più pilotato.}
\label{fig:tx2rx}
\end{figure}

\subsection{Macchina a stati consigliata}
\begin{center}
\begin{tikzpicture}[node distance=12mm and 22mm,font=\sffamily\footnotesize,
   st/.style={draw=sbblue,thick,fill=sblight,rounded corners=3pt,minimum width=26mm,minimum height=9mm,align=center},
   ar/.style={-{Stealth},thick,sbblue2},lb/.style={font=\sffamily\scriptsize,align=center,fill=white,inner sep=1pt}]
 \node[st] (off) {OFF\\SDN = H, PC2 = H};
 \node[st,right=of off] (rdy) {READY\\PC2 = H};
 \node[st,right=of rdy] (rx) {RX\\PC2 = H};
 \node[st,below=of rdy] (arm) {TX ARMATO\\PC2 = L, Si4464 READY};
 \node[st,draw=sbred,fill=sbredbg,below=of arm] (tx) {TX\\PC2 = L};
 \node[st,right=of tx] (end) {FINE TX\\Si4464 READY, PC2 = L};
 \draw[ar] (off) -- node[lb,above]{reset, $t_d$,\\POWER\_UP, config} (rdy);
 \draw[ar] (rdy.20) -- node[lb,above]{START\_RX} (rx.160);
 \draw[ar] (rx.200) -- node[lb,below]{CHANGE\_STATE} (rdy.340);
 \draw[ar] (rdy) -- node[lb,right]{PC2 = L} (arm);
 \draw[ar] (arm) -- node[lb,right]{$\geq\SI{10}{\micro\second}$, START\_TX} (tx);
 \draw[ar] (tx) -- node[lb,above]{PACKET\_SENT} (end);
 \draw[ar] (end) |- node[lb,pos=0.25,right]{$\geq\SI{10}{\micro\second}$, PC2 = H,\\verifica F-10} ($(rdy.east)+(0,-3mm)$);
\end{tikzpicture}
\end{center}
PC2 cambia solo nelle transizioni READY$\leftrightarrow$TX ARMATO e FINE TX$\to$READY, mai negli stati TX.

\section{Limiti di potenza}
\label{sec:limiti}

\subsection{Tabella dei limiti}
{\small
\rowcolors{2}{white}{sbrow}
\begin{xltabular}{\linewidth}{@{}L{42mm}L{30mm}L{30mm}Y@{}}
\caption{Limiti di potenza e tensione della catena TX/RX.}\label{tab:limiti}\\
\tabhead\hd{Punto}&\hd{Limite}&\hd{Valore atteso}&\hd{Fonte e note}\\
\endfirsthead
\tabhead\hd{Punto}&\hd{Limite}&\hd{Valore atteso}&\hd{Fonte e note}\\
\endhead
Uscita del Si4464 (PA\_PWR\_LVL) & \textbf{mai +20 dBm} & circa \SI{0}{\dBm}, da tarare (M-01) & \fonte{LOG F-01; REV-1104 §2.2}\\
Attenuatore Si4464$\to$PA & \textendash & \SI{10.76}{dB} & \fonte{REV-1104 §2.2, calcolo del revisore}\\
Ingresso del PA & \SI{+13}{\dBm} abs.\ max & \SI{-10.8}{\dBm}; \SI{+9.2}{\dBm} con Si4464 a +20 dBm & \fonte{DS-QPA9510 p.\,2; REV-1104 §2.2}\\
Guadagno / compressione del PA & \textendash & $\approx$\SI{42.5}{dB}, P1dB $\approx$\SI{32}{\dBm} a \SI{430}{\mega\hertz} (a \SI{3.6}{\volt}; a \SI{3.3}{\volt} un po' meno) & \fonte{EVB Qorvo pp.\,6--7; REV-1104 §4}\\
Tensione APC & \SI{3.0}{\volt} abs.\ max & \SI{2.87}{\volt} a vuoto (3V3 = \SI{3.3}{\volt}); \SI{2.96}{\volt} max (3V3 = \SI{3.405}{\volt}); \SIrange{2.7}{2.8}{\volt} sotto carico & \fonte{DS-QPA9510 p.\,2; LOG rev.\ 1.45, P09}\\
Filtro FL201 & \SI{2}{\watt} fino a circa \SI{83}{\celsius} (\SI{3.5}{\watt} a \SI{25}{\celsius}) & perdita \SI{1.15}{dB} & \fonte{REV-1104 §2.2}\\
Switch, percorso passante & \SI{33.5}{\dBm} (VDD \SI{5}{\volt}, \SI{85}{\celsius}) & \SIrange{30.3}{30.7}{\dBm}; \SI{32.5}{\dBm} con Si4464 a +20 dBm & \fonte{DS-HMC349A p.\,4; LOG 1.29-01, P-03}\\
Switch, hot switching & \SI{30}{\dBm} (\SI{85}{\celsius}) & \textbf{0}: PC2 si commuta senza RF & \fonte{DS-HMC349A p.\,4; LOG F-02}\\
Switch, porta terminata & \SI{26.5}{\dBm} (\SI{85}{\celsius}) & \textbf{0}: vedi F-14 & \fonte{DS-HMC349A p.\,4; LOG P-02}\\
Uscita d'antenna (MMCX) & \textbf{$\leq$ \SI{31}{\dBm}} & \SIrange{29}{30}{\dBm} & \fonte{LOG F-01, M-06}\\
LNA in TX (perdita dello switch) & \SI{+22}{\dBm} & $\approx$\SI{-14.5}{\dBm} & \fonte{DS-TQP3M9036 p.\,2; REV-1104 §2.2}\\
Ingresso RX del Si4464 & \SI{+10}{\dBm} abs.\ max & dipende dal segnale in antenna & \fonte{DS-Si4464 p.\,13}\\
\bottomrule
\end{xltabular}}

Calcolo della tensione APC a vuoto: $V_\text{APC} = V_\text{3V3}\cdot R_{243}/(R_{242}+R_{243}) = 3{,}3\cdot 220/253 = \SI{2.870}{\volt}$; con $V_\text{3V3} = \SI{3.405}{\volt}$ si ottiene \SI{2.961}{\volt}. La resistenza di Thevenin è \SI{28.7}{\ohm}: con la corrente APC di \SI{4.6}{\milli\ampere} \fonte{DS-QPA9510 p.\,3} la tensione scende di circa \SI{0.13}{\volt}. Il firmware non può regolare l'APC: è fissa. La potenza si regola solo con PA\_PWR\_LVL.

\subsection{Diagramma dei livelli}
\begin{center}
\begin{tikzpicture}
\begin{axis}[width=0.92\linewidth,height=7.2cm,
  symbolic x coords={Si4464,ingresso PA,uscita PA,dopo FL201,MMCX},xtick=data,
  ylabel={Potenza [\si{\dBm}]},ymin=-15,ymax=38,grid=major,grid style={gray!20},
  tick label style={font=\footnotesize\sffamily},label style={font=\small\sffamily},
  legend style={font=\scriptsize\sffamily,at={(0.02,0.98)},anchor=north west},
  enlarge x limits=0.08]
\addplot[very thick,sbgreen,mark=*] coordinates {(Si4464,0) (ingresso PA,-10.8) (uscita PA,31.5) (dopo FL201,30.3) (MMCX,29.5)};
\addlegendentry{obiettivo (Si4464 $\approx$ \SI{0}{\dBm}, stima)}
\addplot[very thick,sbred,mark=square*,dashed] coordinates {(Si4464,20) (ingresso PA,9.2) (uscita PA,33.7) (dopo FL201,32.5) (MMCX,31.7)};
\addlegendentry{Si4464 a +20 dBm (vietato, stima)}
\addplot[thick,sborange,densely dotted,domain=0:4] coordinates {(Si4464,33.5) (MMCX,33.5)};
\addlegendentry{switch passante \SI{33.5}{\dBm}}
\addplot[thick,sbblue,densely dashdotted] coordinates {(Si4464,31) (MMCX,31)};
\addlegendentry{F-01: antenna $\leq$ \SI{31}{\dBm}}
\addplot[thick,sbgray,dashed] coordinates {(Si4464,13) (MMCX,13)};
\addlegendentry{ingresso PA \SI{13}{\dBm} abs.\ max}
\end{axis}
\end{tikzpicture}
\end{center}
{\footnotesize Stime ricavate dai valori della tabella~\ref{tab:limiti}: obiettivo costruito a ritroso da \SI{29.5}{\dBm} all'MMCX (M-06) con \SI{0.8}{dB} di perdita dello switch \fonte{REV-1104 §4} e \SI{1.15}{dB} di FL201; caso a +20 dBm da \SI{32.5}{\dBm} al commutatore \fonte{LOG 1.29-01}. In saturazione la potenza d'uscita resta vicina alla Psat qualunque sia il pilotaggio \fonte{LOG R-04}.}

\begin{regola}[F-01: limite rigido di potenza]
\begin{itemize}
\item PA\_PWR\_LVL del Si4464 è una \textbf{costante} tarata sul prototipo con la misura M-01, per avere \textbf{$\leq$ \SI{31}{\dBm}} all'uscita d'antenna. Non deve essere modificabile da telecomando; al massimo il firmware satura il valore richiesto \fonte{LOG F-01}.
\item Il Si4464 non deve mai trasmettere a +20 dBm.
\item Scrivere PA\_PWR\_LVL dopo ogni POWER\_UP e dopo ogni caricamento della configurazione WDS, prima del primo START\_TX. Secondo la documentazione API il valore di default è \code{0x7F}, cioè il massimo dei 128 passi: \DV{} sulla API della revisione B1.
\item Se M-05 (armoniche) non passa, ridurre la potenza di \SIrange{2}{3}{dB} \fonte{LOG M-05}.
\end{itemize}
\end{regola}

\begin{regola}[F-03: durata della trasmissione]
Preferire trasmissioni a impulsi o con duty cycle limitato. L242 e FL201 scaldano in trasmissione continua alla massima potenza. Se serve un TX continuo, ridurre il valore F-01 \fonte{LOG F-03}. La temperatura del PA si legge su PB1.
\end{regola}

\begin{attenzione}[carico in uscita]
Non trasmettere mai senza antenna o senza un carico da \SI{50}{\ohm} adatto alla potenza sull'MMCX. Il datasheet del Si4464 avverte che il PA si può danneggiare se acceso senza carico \fonte{DS-Si4464 p.\,13}; il QPA9510 garantisce l'assenza di danni fino a VSWR 10:1 solo nelle condizioni di prova del datasheet \fonte{DS-QPA9510 p.\,3}.
\end{attenzione}

\section{TCXO e configurazione base del Si4464}
\label{sec:si-config}

\subsection{Clock di riferimento}
\begin{itemize}
\item TCXO Y201 ECS-TXO-20CSMV4-320-AY: \SI{32.000}{\mega\hertz}, onda sinusoidale troncata, stabilità codice A = \SI{\pm0.5}{\ppm}, temperatura codice Y = \SIrange{-30}{85}{\celsius} \fonte{DS-ECS pp.\,1--2}.
\item Accoppiato in AC a XIN tramite C204 \SI{10}{\nano\farad} e R201 \SI{33}{\ohm}; XOUT è flottante; test point TP701 su XIN \fonte{SCH}. Il Si4464 consiglia \SIrange{0.6}{1.4}{\volt} picco-picco in AC, massimo \SI{1.8}{\volt} \fonte{DS-Si4464 p.\,38}.
\item Errore di frequenza dovuto alla stabilità: $435\cdot10^6\times0{,}5\cdot10^{-6} = \SI{\pm217.5}{\hertz}$ (calcolo). La tolleranza iniziale a \SI{25}{\celsius} va letta sul datasheet ECS: \DV.
\end{itemize}

\begin{verifica}[campo di temperatura del TCXO]
La versione \code{-AY} è specificata da \SI{-30}{\celsius}; per \SI{-40}{\celsius} servirebbe la \code{-AN} \fonte{REV-1104 P07; LOG R-07}. Sotto \SI{-30}{\celsius} la frequenza non è garantita.
\end{verifica}

\subsection{Configurazione}
\begin{regola}[F-04: clock del Si4464]
POWER\_UP con l'opzione TCXO (TCXO a \SI{32}{\mega\hertz} su XIN, XOUT flottante) e banco di capacità del quarzo GLOBAL\_XO\_TUNE = 0 \fonte{LOG F-04; DS-Si4464 p.\,38}.
\end{regola}

\begin{xltabular}{\linewidth}{@{}L{40mm}L{40mm}Y@{}}
\caption{Parametri e comandi principali. I codici vengono dalla documentazione API Silicon Labs consultata in rete (revisione C2A); per la Si4464-B1B vanno confermati con WDS: \DV.}\\
\tabhead\hd{Elemento}&\hd{Valore}&\hd{Note}\\
\endfirsthead
\tabhead\hd{Elemento}&\hd{Valore}&\hd{Note}\\
\endhead
Frequenza portante & \SI{435}{\mega\hertz} & banda Si4464 \SIrange{353}{479}{\mega\hertz} \fonte{DS-Si4464 p.\,9}\\
\code{POWER\_UP} & comando \code{0x02} & BOOT\_OPTIONS: modo PRO; XTAL\_OPTIONS: bit TCXO = 1\\
XO\_FREQ & \num{32000000} = \code{0x01E84800} & 4 byte, MSB per primo (conversione decimale/esadecimale: calcolo)\\
\code{GLOBAL\_XO\_TUNE} & 0 (gruppo \code{0x00}, indice \code{0x00}) & default \code{0x40}: va scritto a 0 (F-04)\\
\code{PA\_PWR\_LVL} & gruppo \code{0x22}, indice \code{0x01} & 128 passi \fonte{DS-Si4464 p.\,35}; valore = costante F-01\\
\code{PART\_INFO} & comando \code{0x01} & risposta PART = \code{0x4464} attesa\\
\code{GET\_INT\_STATUS} & comando \code{0x20} & nella ISR di nIRQ\\
\code{START\_TX} / \code{START\_RX} & \code{0x31} / \code{0x32} & stato finale configurabile (READY)\\
\code{CHANGE\_STATE} & \code{0x34} & per forzare READY\\
\code{READ\_CMD\_BUFF} & \code{0x44} & CTS = \code{0xFF} \fonte{DS-Si4464 p.\,17}\\
\bottomrule
\end{xltabular}

\subsection{Procedura di inizializzazione della radio}
\begin{enumerate}
\item PC2 alto (RX, PA spento) e NSEL (PA8) alto.
\item Modo A: PC4 alto; attendere almeno \SI{300}{\milli\second} ($t_d$ max). Modo B: PC4 alto per almeno \SI{10}{\micro\second}, poi basso.
\item Polling di CTS (\code{0xFF}) con timeout.
\item \code{POWER\_UP} con opzione TCXO e XO\_FREQ = \SI{32}{\mega\hertz}; attendere CTS.
\item \code{PART\_INFO}: controllare che il chip sia un Si4464. In caso contrario segnalare l'errore in telemetria e non trasmettere.
\item Caricare la configurazione generata da WDS (\code{radio\_config.h}: Si4464, TCXO \SI{32}{\mega\hertz}, \SI{435}{\mega\hertz}, modulazione e data rate della missione, PA a +20 dBm classe E come da match AN648 \fonte{SCH nota UHF}).
\item Scrivere GLOBAL\_XO\_TUNE = 0 e PA\_PWR\_LVL = costante F-01 (sovrascrivono i valori di WDS).
\item Configurare nIRQ (interrupt PACKET\_SENT, PACKET\_RX, errori CRC) e passare in RX o READY.
\end{enumerate}

\section{Protezione della ricezione}
\label{sec:rx-protezione}
Il massimo assoluto all'ingresso RX del Si4464 è \SI{+10}{\dBm} \fonte{DS-Si4464 p.\,13}. Davanti c'è un LNA da circa \SI{21}{dB} \fonte{REV-1104 §4}. Il segnale massimo ammesso all'MMCX è quindi circa $10 - 21 + 0{,}8 + 0{,}17 \approx \SI{-10}{\dBm}$ (calcolo: guadagno LNA, perdite di switch e filtro RX \fonte{REV-1104 §4, P11}).

\begin{attenzione}[al banco]
\begin{itemize}
\item Non collegare mai un generatore o un altro trasmettitore direttamente all'MMCX senza attenuatori. Tenere il segnale all'MMCX sotto \SI{-20}{\dBm} (\SI{10}{dB} di margine sul calcolo sopra).
\item Prova tra due schede SYMPHONY: con \SI{30}{\dBm} in uscita servono almeno \SI{50}{dB} di attenuazione ($30 - 50 = \SI{-20}{\dBm}$; calcolo).
\item La revisione indica che l'abs max si supera già con segnali in antenna oltre circa \SI{-5}{\dBm} \fonte{REV-1104 §2.2}.
\end{itemize}
\end{attenzione}

\section{Supervisore e verifica dopo la trasmissione}
\label{sec:rf-supervisore}
In TX il PA assorbe circa \SI{1.7}{\ampere} dal 3V3. Se il rail scende sotto la soglia di IC205 (\SI{2.88}{\volt} min, \SI{3.00}{\volt} max), il supervisore spegne il Si4464 per circa $t_d$ e la trasmissione si interrompe \fonte{LOG §46}. Il danno non c'è, ma il Si4464 perde la configurazione.

\begin{regola}[F-10: controllo dopo ogni TX]
Dopo ogni trasmissione controllare che il Si4464 risponda ancora (CTS, PART\_INFO o stato del chip). Se è stato resettato: reinizializzarlo (procedura del paragrafo~\ref{sec:si-config}), contare l'evento e segnalarlo in telemetria \fonte{LOG F-10}. Utile anche registrare VDDA (VREFINT) subito dopo la TX.
\end{regola}

\section{Collaudo RF}
\label{sec:collaudo}
Il firmware di collaudo deve permettere di: trasmettere a PA\_PWR\_LVL fisso e noto, ripetere la trasmissione a impulsi o in continuo per un tempo limitato, ricevere in continuo, leggere temperature e VDDA dalla console.

{\small
\rowcolors{2}{white}{sbrow}
\begin{xltabular}{\linewidth}{@{}L{10mm}L{42mm}L{40mm}Y@{}}
\caption{Misure di collaudo collegate al firmware \fonte{LOG M-01\ldots M-06}.}\label{tab:collaudo}\\
\tabhead\hd{ID}&\hd{Misura}&\hd{Criterio}&\hd{Cosa serve dal firmware}\\
\endfirsthead
\tabhead\hd{ID}&\hd{Misura}&\hd{Criterio}&\hd{Cosa serve dal firmware}\\
\endhead
M-01 & Pout all'uscita d'antenna su \SI{50}{\ohm} (attenuatore di potenza), a temperatura ambiente e se possibile a caldo, con il Si4464 al valore F-01 e al massimo & $\leq$ \SI{31}{\dBm} con il valore F-01; altrimenti ridurre la VAPC con R242/R243 & build di collaudo con PA\_PWR\_LVL impostabile da console; fissa il valore definitivo di F-01\\
M-02 & Adattamento d'uscita del PA (S22 o potenza in banda \SIrange{430}{440}{\mega\hertz}) & \textendash & TX a frequenza programmabile\\
M-03 & 3V3 al pin 5 di IC205 durante una TX a piena potenza (valore F-01) & minimo $\geq$ \SI{3.05}{\volt}; altrimenti modo B e aggiornamento di F-05 & TX ripetuta; costante \code{MODO\_A/B}\\
M-04 & NRST con R413 DNP e firmware bloccato & impulso basso \SIrange{140}{280}{\milli\second} ogni \SIrange{1.12}{2.24}{\second}; con R413 montata nessun reset & comando che blocca il refresh di WDI\\
M-05 & Armoniche H2 e H3 al connettore d'antenna, TX alla potenza massima & entro la specifica di missione; altrimenti \SIrange{-2}{-3}{dB} via firmware & TX continua limitata nel tempo\\
M-06 & Potenza all'MMCX e corrente del PA a pilotaggio fisso & attesa \SIrange{29}{30}{\dBm}; se più bassa, accordare C260/C262 e ripetere M-05 & PA\_PWR\_LVL fisso\\
\bottomrule
\end{xltabular}}

\begin{nota}[set-up consigliato]
MMCX $\to$ attenuatore di potenza da \SI{50}{\ohm} dimensionato per almeno \SI{2}{\watt} continui $\to$ misuratore di potenza o analizzatore di spettro. Sonde dell'oscilloscopio: TP702 (PC2), TP703 (EN\_RX\_LNA), TP701 (XIN del Si4464) per verificare le sequenze dei paragrafi~\ref{sec:sequenze}. Le misure M-01, M-03 e M-06 vanno ripetute con il firmware definitivo.
\end{nota}

\chapter{Supervisione e affidabilità}
\label{ch:supervisione}

\section{Watchdog esterno STWD100 (U403)}
\subsection{Collegamenti}
\begin{xltabular}{\linewidth}{@{}L{20mm}L{34mm}Y@{}}
\tabhead\hd{Pin U403}&\hd{Collegato a}&\hd{Note}\\
1 \(\overline{\mathrm{WDO}}\) & NRST dell'MCU & Open-drain, attivo basso; pull-up R2 \SI{10}{\kilo\ohm}\\
3 \(\overline{\mathrm{EN}}\) & R413 \SI{0}{\ohm} \textbf{DNP} verso 3V3 & Pull-down interno: EN flottante = watchdog \textbf{attivo}\\
4 WDI & PC6 (\net{/MCU/WDT\_WDI}) & Refresh dal firmware\\
5 VCC & 3V3 (C425 \SI{100}{\nano\farad}) & \\
\bottomrule
\end{xltabular}
\fonte{LOG rev.\ 1.41; SCH}

\subsection{Tempi (versione NY)}
\begin{xltabular}{\linewidth}{@{}L{46mm}C{18mm}C{18mm}C{18mm}Y@{}}
\tabhead\hd{Parametro}&\hd{min}&\hd{tip}&\hd{max}&\hd{Fonte}\\
Timeout $t_{WD}$ & \SI{1.12}{\second} & \SI{1.6}{\second} & \SI{2.24}{\second} & \fonte{LOG rev.\ 1.41 (DS Rev 10, tab.\,4)}\\
Impulso di reset $t_{PW}$ & \SI{140}{\milli\second} & \SI{210}{\milli\second} & \SI{280}{\milli\second} & idem\\
Pull-down interno su EN & \SI{32}{\kilo\ohm} & \SI{63}{\kilo\ohm} & \SI{100}{\kilo\ohm} & idem\\
Impulso minimo su WDI ed EN & \SI{1}{\micro\second} & & & idem; glitch $<$\SI{100}{\nano\second} ignorati\\
\bottomrule
\end{xltabular}
Nella versione Y il timer si azzera sia sul fronte di salita sia su quello di discesa di WDI \fonte{DS-STWD100 Rev 8 p.\,14 (pagina stampata 11)}.

\begin{regola}[F-11, F-12, F-13: watchdog]
\begin{itemize}
\item \textbf{F-11}: commutare PC6 almeno ogni circa \SI{1}{\second}, anche durante le operazioni lunghe ($t_{WD}$ min \SI{1.12}{\second}). Il fronte va generato \textbf{dal software nel ciclo principale}, dopo aver verificato che il sistema funziona. Mai con un timer in PWM o con un interrupt a tempo autonomo: il timer continuerebbe a girare con la CPU bloccata \fonte{LOG F-11, rev.\ 1.42}.
\item \textbf{F-12}: primo fronte entro \SI{1.12}{\second} dall'accensione, prima delle inizializzazioni lente (per esempio l'attesa di \SI{300}{\milli\second} del Si4464) \fonte{LOG F-12}.
\item \textbf{F-13}: all'avvio leggere \code{RCC\_CSR} per la causa del reset (pin NRST = watchdog esterno o debugger) e riportarla in telemetria \fonte{LOG F-13}.
\end{itemize}
\end{regola}

\begin{nota}[schema di refresh consigliato]
Un approccio robusto è il \enquote{watchdog a richiesta}: ogni task critico alza un suo flag; il ciclo principale commuta PC6 solo se tutti i flag sono alzati, poi li azzera. Un periodo di commutazione di qualche centinaio di millisecondi lascia margine rispetto a \SI{1.12}{\second}. Nelle attese lunghe (scritture, sequenze di avvio) inserire refresh espliciti.
\end{nota}

\begin{attenzione}[R413 e configurazione di volo]
R413 non montata = watchdog sempre attivo (volo). R413 \SI{0}{\ohm} montata = EN alto, watchdog spento (debug SWD, bootloader UART) \fonte{SCH nota MCU; LOG rev.\ 1.42}. Un GPIO su EN è stato scartato di proposito: un firmware fuori controllo non deve poter disattivare il proprio watchdog. R413 non ha il flag \enquote{Exclude from BOM}: il controllo della BOM prima dell'ordine è procedurale \fonte{LOG P-06}.
\end{attenzione}

Dopo un reset dell'MCU PC6 è in alta impedenza. Se il firmware non riparte entro \SI{1.12}{\second}, il watchdog resetta di nuovo: è il comportamento voluto \fonte{LOG rev.\ 1.41}.

\section{Supervisore TPS3825 (IC205) e RF\_MR\_N}
\begin{itemize}
\item Soglia sul 3V3: \SI{2.88}{}/\SI{2.93}{}/\SI{3.00}{\volt} (min/tip/max) \fonte{LOG §46}.
\item Ritardo dopo il rilascio: $t_d$ = \SI{120}{}/\SI{200}{}/\SI{300}{\milli\second} \fonte{DS-TPS3825 p.\,8}.
\item \(\overline{\mathrm{MR}}\): impulso minimo \SI{1}{\micro\second} \fonte{DS-TPS3825 p.\,8}; soglie \(0{,}3\,V_\text{DD}\) e \(0{,}7\,V_\text{DD}\); pull-up interno \SI{90}{\kilo\ohm} \fonte{DS-TPS3825 pp.\,5, 7} più R208 \SI{10}{\kilo\ohm} esterno.
\item Uscita RESET (attiva alta) $\to$ R204 $\to$ SDN del Si4464. Uscita \(\overline{\mathrm{RESET}}\) $\to$ R209 \SI{1.5}{\kilo\ohm} $\to$ LED D203 \fonte{SCH}.
\item Uso da firmware e modi A/B: paragrafo~\ref{sec:sdn}. Misura M-03: paragrafo~\ref{sec:collaudo}.
\end{itemize}

\section{RTC DS3231M (U404)}
\subsection{Collegamenti}
I2C2 (PB10/PB11, pull-up \SI{10}{\kilo\ohm}), indirizzo \code{0x68}. \(\overline{\mathrm{RST}}\) su PB9. INT/SQW ha solo il pull-up R411 e non arriva all'MCU; l'uscita 32\,kHz è scollegata (R410 rimossa) \fonte{SCH; LOG rev.\ 1.45, P03}. Alimentazione di backup da 10 supercondensatori CPH3225A (\SI{110}{\milli\farad}) caricati dal 3V3 tramite BAT54WS e \SI{1}{\kilo\ohm}: il requisito è tenere l'ora per pochi secondi; l'autonomia stimata è di circa \SI{7}{\hour} \fonte{LOG R-06, R-24}.

\subsection{Registri da impostare all'avvio}
\begin{xltabular}{\linewidth}{@{}L{18mm}L{36mm}L{28mm}Y@{}}
\caption{Registri di controllo del DS3231M \fonte{DS-DS3231M pp.\,11, 13, 15}.}\\
\tabhead\hd{Registro}&\hd{Bit}&\hd{Valore al POR}&\hd{Impostazione consigliata}\\
\endfirsthead
\tabhead\hd{Registro}&\hd{Bit}&\hd{Valore al POR}&\hd{Impostazione consigliata}\\
\endhead
\code{0Eh} Control & EOSC (7), BBSQW (6), CONV (5), INTCN (2), A2IE (1), A1IE (0) & EOSC = 0, BBSQW = 0, INTCN = 1, A1IE = A2IE = 0 & Lasciare EOSC = 0 (oscillatore attivo anche in batteria), BBSQW = 0, INTCN = 1, allarmi disabilitati: INT/SQW non è collegato.\\
\code{0Fh} Status & OSF (7), EN32KHZ (3), BSY (2), A2F (1), A1F (0) & OSF = 1, EN32KHZ = 1 & \textbf{EN32KHZ = 0} (F-16). OSF: leggerlo, poi azzerarlo dopo aver scritto un'ora valida. Azzerare A1F e A2F.\\
\code{10h} Aging & SIGN + DATA & 0 & Lasciare 0 (non serve per la precisione dichiarata) \fonte{DS-DS3231M p.\,15}.\\
\code{00h}--\code{06h} & ora e data (BCD) & non definiti al primo POR & Scrivere l'ora (da terra o dal dock) quando OSF = 1.\\
\bottomrule
\end{xltabular}

\begin{regola}[F-09 e F-16: DS3231M]
All'avvio scrivere EN32KHZ = 0 (registro \code{0Fh}, bit 3): il pin 32\,kHz non è usato e l'uscita abilitata consuma la carica di backup (circa \SI{0.85}{\micro\ampere} contro \SI{0.75}{\micro\ampere} a \SI{3}{\volt}) \fonte{LOG F-16}. Dopo un'interruzione lunga controllare il flag OSF e risincronizzare l'ora \fonte{LOG F-09}.
\end{regola}

\begin{nota}[F-09 e F-16]
F-09 (rev.\ 1.29) e F-16 (rev.\ 1.45) contengono la stessa regola su EN32KHZ. F-09 aggiunge il controllo di OSF.
\end{nota}

\subsection{Pin RST (PB9)}
\(\overline{\mathrm{RST}}\) è un'uscita open-drain con pull-up interno da \SI{50}{\kilo\ohm} e anche un ingresso per pulsante con antirimbalzo \fonte{DS-DS3231M p.\,7}. Va basso quando VCC scende sotto $V_\text{PF}$ e resta basso per $t_\text{REC}$ (\SI{250}{\milli\second} tip., \SI{300}{\milli\second} max) dopo il ritorno di VCC \fonte{DS-DS3231M pp.\,3, 10}. Lo stato di RST non influisce su I2C e RTC \fonte{DS-DS3231M p.\,4}.
\begin{itemize}
\item Configurare PB9 come \textbf{ingresso senza pull}. Si può usare come indicatore di power-fail.
\item Non configurare PB9 come uscita push-pull: entrerebbe in conflitto con l'uscita del DS3231M. Portare RST basso da firmware avvia un ciclo di reset del DS3231M \fonte{DS-DS3231M p.\,10}: non serve.
\end{itemize}

\section{FRAM FM25V10 (U401)}
\begin{itemize}
\item \SI{1}{Mbit} (\num{131072} byte), \num{e14} cicli di lettura/scrittura, scrittura senza attesa \fonte{DS-FM25V10 p.\,1}.
\item La FRAM all'accensione ha la scrittura disabilitata: WREN prima di ogni WRITE/WRSR \fonte{DS-FM25V10 p.\,6}.
\item Controllare RDID all'avvio (\code{7F 7F 7F 7F 7F 7F C2 24 00}) come test di presenza del bus SPI2 \fonte{DS-FM25V10 p.\,11}.
\item /WP è fisso alto: la protezione hardware dello status register non c'è. Usare i bit BP1/BP0 se serve proteggere zone della memoria \fonte{DS-FM25V10 p.\,7}.
\item Consiglio di progetto firmware: proteggere i dati critici (offset di calibrazione F-08, contatori di reset F-10/F-13, configurazione) con CRC e doppia copia.
\end{itemize}

\chapter{LED e debug}
\label{ch:led}

\section{LED dell'Engineering Model}
Nessun LED è collegato a un GPIO dell'MCU \fonte{SCH}. I due LED danno un riscontro visivo indipendente dal firmware.

\begin{xltabular}{\linewidth}{@{}L{14mm}L{30mm}L{38mm}Y@{}}
\tabhead\hd{LED}&\hd{Pilotato da}&\hd{Acceso quando}&\hd{Note}\\
D203 (rosso) & \(\overline{\mathrm{RESET}}\) di IC205, tramite R209 \SI{1.5}{\kilo\ohm} & il Si4464 è fuori reset (SDN basso, modo A) & \SI{0.93}{\milli\ampere} continui. Spento durante il reset della radio o con il 3V3 sotto soglia. Nel modo B non indica più lo stato della radio \fonte{LOG rev.\ 1.45, P14; §46}.\\
D9 (verde) & 3V3 tramite R49 \SI{10}{\kilo\ohm} e JP1 & 3V3 presente \textbf{e} JP1 chiuso & JP1 aperto di serie. Luce molto debole (circa \SI{0.09}{\milli\ampere}), voluto: serve solo per il debug iniziale \fonte{LOG S-13}.\\
\bottomrule
\end{xltabular}

\begin{nota}[come usare D203 al banco]
D203 è il modo più rapido per vedere un reset della radio durante la trasmissione (calo del 3V3, misura M-03): se si spegne per circa \SI{200}{\milli\second} durante una TX, IC205 è intervenuto. In parallelo il firmware deve rilevare l'evento con F-10.
\end{nota}

\section{Flight Model}
\begin{regola}[LED sul Flight Model]
Sul Flight Model i LED di debug non si montano: D203 e R209 (e D9) vanno esclusi dalla BOM di volo \fonte{LOG rev.\ 1.45, P14}. Il firmware non deve dipendere dai LED. Ogni informazione utile che oggi si legge dai LED va anche in telemetria (stato della radio, reset, eventi F-10).
\end{regola}

Altre differenze tra configurazione di sviluppo e di volo:
\begin{xltabular}{\linewidth}{@{}L{30mm}L{46mm}Y@{}}
\tabhead\hd{Elemento}&\hd{Sviluppo (banco)}&\hd{Volo}\\
R413 (EN watchdog) & montata: watchdog spento & \textbf{non montata}: watchdog attivo\\
R205 / R204 & modo A di serie; modo B solo se M-03 fallisce & secondo l'esito di M-03 (costante \code{MODO\_A/B})\\
LED D203, D9 & montati & non montati\\
JP201, JP6 & secondo il banco & secondo l'architettura del bus del satellite\\
Firmware & build di collaudo con comandi di test & build di volo: PA\_PWR\_LVL fisso (F-01), nessun comando di test RF\\
\bottomrule
\end{xltabular}

\section{Strumenti di debug}
\begin{itemize}
\item \textbf{SWD} su J1 (TC2030): SWDIO, SWCLK, NRST, SWO, 3V3 e GND. Serve R413 montata per lavorare con i breakpoint.
\item \textbf{Console} su USART1 tramite CH340E e USB-C. La scheda deve avere il 3V3 (EN\_DOCK).
\item \textbf{Bootloader di sistema}: TP1 (BOOT0) alto al reset, STM32CubeProgrammer in UART \fonte{LOG F-07}.
\item \textbf{Test point}: TP701 (XIN Si4464, clock del TCXO), TP702 (PC2), TP703 (EN\_RX\_LNA), TP1 (BOOT0) \fonte{SCH}.
\end{itemize}

\chapter{Note firmware e checklist}
\label{ch:notefw}

\section{Elenco consolidato F-01\ldots F-16}
Le note F-xx sono le richieste che il progetto hardware fa al firmware, raccolte nel log. Colonne della tabella: \textbf{ID} = codice della nota; \textbf{Regola} = che cosa deve fare il firmware; \textbf{Motivo} = che cosa succede all'hardware se non lo fa; \textbf{Fonte} = dove si trova la nota nel log; \textbf{Priorità} = gravità delle conseguenze.

Priorità: \textbf{\textcolor{sbred}{Critica}} = rischio di danno all'hardware o di perdita del satellite; \textbf{\textcolor{sborange}{Alta}} = malfunzionamento sicuro se non rispettata; \textbf{\textcolor{sbgreen}{Media}} = qualità, consumi, diagnostica.

{\small
\rowcolors{2}{white}{sbrow}
\begin{xltabular}{\linewidth}{@{}L{11mm}L{56mm}Y L{20mm}L{14mm}@{}}
\caption{Note per il firmwarista dal log hardware.}\label{tab:notefw}\\
\tabhead\hd{ID}&\hd{Regola}&\hd{Motivo}&\hd{Fonte}&\hd{Priorità}\\
\endfirsthead
\tabhead\hd{ID}&\hd{Regola}&\hd{Motivo}&\hd{Fonte}&\hd{Priorità}\\
\endhead
\multicolumn{5}{r}{\footnotesize\itshape continua nella pagina seguente}\\
\endfoot
\bottomrule
\endlastfoot
F-01 & PA\_PWR\_LVL del Si4464 è una costante tarata con M-01 per avere $\leq$\SI{31}{\dBm} all'antenna; non modificabile da telecomando; mai +20 dBm. & Con il Si4464 al massimo il PA satura e lo switch riceve $\approx$\SI{32.5}{\dBm} contro \SI{33.5}{\dBm} abs.\ max. & LOG R-04, rev.\ 1.29 & \textcolor{sbred}{Critica}\\
F-02 & Entrata in TX: Si4464 non in TX $\to$ PC2 = L $\to$ attesa $\to$ START\_TX. Uscita: Si4464 fuori da TX $\to$ solo dopo PC2 = H. Mai cambiare PC2 mentre il Si4464 trasmette. & Hot switching dell'HMC349A, limite \SI{30}{\dBm}. & LOG R-04, rev.\ 1.29 & \textcolor{sbred}{Critica}\\
F-03 & Trasmissioni a impulsi o con duty cycle limitato; se serve TX continuo ridurre F-01. & L242 e FL201 scaldano in TX continuo. & LOG rev.\ 1.25, R-04 & \textcolor{sborange}{Alta}\\
F-04 & POWER\_UP con opzione TCXO (\SI{32}{\mega\hertz} su XIN) e GLOBAL\_XO\_TUNE = 0. & Clock esterno: il banco di capacità va azzerato. & LOG §2; DS-Si4464 p.\,38 & \textcolor{sborange}{Alta}\\
F-05 & Costante MODO\_A / MODO\_B. Modo A: PC4 L = reset; dopo il rilascio attendere il ritardo di IC205 (fino a \SI{300}{\milli\second}) e il POR, poi CTS. Modo B: logica invertita, SDN alto $\geq$\SI{10}{\micro\second}. & Supervisore IC205 escludibile (R204/R205). & LOG R-19, rev.\ 1.34 & \textcolor{sborange}{Alta}\\
F-06 & HSE\_VALUE = 24000000. & Quarzo Y1 da \SI{24}{\mega\hertz}: altrimenti UART, CAN e timer sbagliati. & LOG S-04, rev.\ 1.15 & \textcolor{sborange}{Alta}\\
F-07 & Console su USART1 PA9/PA10 (AF7) verso il CH340E, non UART4. Bootloader: BOOT0 (TP1) alto al reset, STM32CubeProgrammer in UART. & Collegamento hardware della console. & LOG rev.\ 1.8 & \textcolor{sbgreen}{Media}\\
F-08 & ADC1 canale 16 (PB1, PA) e canale 2 (PC1, buck), misura raziometrica su VDDA; calibrazione a un punto per canale. & Offset del TLV9051 $\approx$\SI{\pm0.9}{\celsius}. & LOG S-06, R-18 & \textcolor{sbgreen}{Media}\\
F-09 & All'avvio EN32KHZ = 0 (\code{0Fh} bit 3); dopo un'interruzione lunga controllare OSF e risincronizzare l'ora. & Uscita 32\,kHz non usata; ora non valida dopo lo stop dell'oscillatore. & LOG S-09, rev.\ 1.20 & \textcolor{sbgreen}{Media}\\
F-10 & Dopo ogni TX controllare che il Si4464 risponda (CTS, PART\_INFO, stato); se resettato reinizializzare, contare, segnalare in telemetria. & Possibile reset da calo del 3V3 in TX (IC205). & LOG R-19, rev.\ 1.34 & \textcolor{sborange}{Alta}\\
F-11 & Commutare PC6 (WDI) almeno ogni $\approx$\SI{1}{\second}, dal ciclo principale; mai con timer/PWM o interrupt autonomo. & $t_{WD}$ min \SI{1.12}{\second}; un timer girerebbe anche con la CPU bloccata. & LOG rev.\ 1.41--1.42 & \textcolor{sbred}{Critica}\\
F-12 & Primo impulso su WDI entro \SI{1.12}{\second} dall'accensione, prima delle inizializzazioni lente. & Altrimenti reset ciclico all'avvio. & LOG rev.\ 1.41 & \textcolor{sbred}{Critica}\\
F-13 & All'avvio leggere RCC\_CSR (causa del reset) e riportarla in telemetria. & Diagnosi dei reset del watchdog esterno. & LOG rev.\ 1.41 & \textcolor{sbgreen}{Media}\\
F-14 & TX$\to$RX: prima fermare la TX (Si4464 fuori da TX), attendere $\geq$\SI{10}{\micro\second}, solo dopo PC2 verso RX. Mai commutare PC2 durante o subito dopo la TX. & Porta terminata dell'HMC349A: \SI{26.5}{\dBm} abs.\ max; l'APC si scarica in alcuni \si{\micro\second}. & LOG P-02, rev.\ 1.43 & \textcolor{sbred}{Critica}\\
F-15 & All'avvio PB12 (CS FRAM) e PA8 (NSEL Si4464) uscite alte prima di ogni inizializzazione SPI. & Nessun pull-up esterno sui CS. & LOG P-12, rev.\ 1.44 & \textcolor{sbred}{Critica}\\
F-16 & All'avvio EN32kHz = 0 nel DS3231M. & Risparmio sulla carica di backup. & LOG P03, rev.\ 1.45 & \textcolor{sbgreen}{Media}\\
\end{xltabular}}

\section{Checklist di avvio del firmware}
L'ordine è importante: le prime voci proteggono l'hardware e il watchdog.

\newcommand{\chk}{\item[$\square$]}
\begin{enumerate}[label=\arabic*.,leftmargin=*]
\item \textbf{Prime istruzioni dopo il reset} (prima di \code{SystemInit} lunghe o di copie di memoria lente)
 \begin{itemize}[leftmargin=8mm]
 \chk PA8 (NSEL) e PB12 (CS FRAM) uscite push-pull a livello alto (F-15).
 \chk PC2 uscita push-pull a livello \textbf{alto}: switch su RX, PA spento.
 \chk PC4 uscita a livello alto in entrambi i modi: nel modo A la radio esce dal reset dopo $t_d$, nel modo B resta spenta.
 \chk PC6 uscita e primo fronte su WDI (F-12). Annotare l'istante.
 \end{itemize}
\item \textbf{Diagnostica del reset}
 \begin{itemize}[leftmargin=8mm]
 \chk Leggere e salvare \code{RCC\_CSR}, poi azzerare i flag (F-13).
 \end{itemize}
\item \textbf{Clock}
 \begin{itemize}[leftmargin=8mm]
 \chk HSE \SI{24}{\mega\hertz} (\code{HSE\_VALUE = 24000000}) e PLL (F-06). Gestire il guasto dell'HSE (CSS) con un ripiego su MSI/HSI.
 \end{itemize}
\item \textbf{GPIO}
 \begin{itemize}[leftmargin=8mm]
 \chk Pin liberi in modo analogico (tabella del paragrafo~\ref{sec:pin-liberi}).
 \chk PC0 ingresso con pull-up, EXTI0 sul fronte di discesa (non ancora abilitato).
 \chk PB9 ingresso senza pull. PC13--PC15 ingressi (in attesa di decisione, §\ref{sec:pc13}).
 \end{itemize}
\item \textbf{Periferiche}
 \begin{itemize}[leftmargin=8mm]
 \chk USART1 (console), FDCAN1, ADC (canali 2 e 16, VREFINT).
 \chk SPI2 e FRAM: RDID, lettura della configurazione e degli offset di calibrazione (F-08).
 \chk I2C2 e DS3231M: EN32KHZ = 0 (F-16), controllo di OSF (F-09), \code{0Eh} = INTCN 1, allarmi off.
 \end{itemize}
\item \textbf{Radio} (refresh di WDI durante le attese)
 \begin{itemize}[leftmargin=8mm]
 \chk Attendere almeno \SI{300}{\milli\second} dal rilascio di PC4 (modo A), polling di CTS con timeout (F-05).
 \chk SPI1 modo 0, $\leq$\SI{10}{\mega\hertz}.
 \chk POWER\_UP con TCXO, XO\_FREQ \SI{32}{\mega\hertz}; GLOBAL\_XO\_TUNE = 0 (F-04).
 \chk PART\_INFO; configurazione WDS; \textbf{PA\_PWR\_LVL = costante F-01} (F-01).
 \chk Abilitare EXTI0 (nIRQ) e passare in RX o READY.
 \end{itemize}
\item \textbf{Ciclo principale}
 \begin{itemize}[leftmargin=8mm]
 \chk Refresh di WDI condizionato allo stato dei task (F-11).
 \chk Telemetria: causa del reset, contatore F-10, temperature, VDDA.
 \end{itemize}
\end{enumerate}

\section{Checklist per ogni trasmissione}
\begin{itemize}[leftmargin=8mm]
\chk Antenna o carico da \SI{50}{\ohm} presenti (al banco).
\chk Temperatura del PA (PB1) entro i limiti scelti dalla missione; duty cycle limitato (F-03).
\chk Si4464 in READY, PA\_PWR\_LVL = F-01.
\chk PC2 = L, attesa $\geq$\SI{10}{\micro\second}, START\_TX (F-02).
\chk Fine pacchetto, Si4464 fuori da TX, attesa $\geq$\SI{10}{\micro\second}, PC2 = H (F-14).
\chk Verifica che il Si4464 risponda; se no reinizializzare e contare l'evento (F-10).
\end{itemize}

\section{Errori da non fare}
\begin{attenzione}[riepilogo]
\begin{itemize}
\item Commutare PC2 mentre il Si4464 trasmette o subito dopo (F-02, F-14).
\item Lasciare PA\_PWR\_LVL al valore di default o renderlo modificabile da terra (F-01).
\item Generare WDI con un timer o con DMA (F-11).
\item Inizializzare SPI1 o SPI2 prima di aver portato alti PA8 e PB12 (F-15).
\item Configurare PB9 (RST del DS3231M) come uscita push-pull.
\item Collegare un generatore all'MMCX senza attenuatori (paragrafo~\ref{sec:rx-protezione}).
\item Tenere PC2 basso fuori dalla trasmissione: il PA assorbe la corrente di riposo e scalda.
\end{itemize}
\end{attenzione}

\appendix
\chapter{Componenti e datasheet}
\label{app:componenti}

Le pagine sono quelle del PDF del datasheet usato per questo manuale. Le sigle tra parentesi quadre nel testo (DS-xxx) rimandano a questa tabella.

{\small
\rowcolors{2}{white}{sbrow}
\begin{xltabular}{\linewidth}{@{}L{17mm}L{33mm}L{36mm}Y@{}}
\tabhead\hd{Rif.}&\hd{Parte}&\hd{Datasheet (sigla)}&\hd{Pagine chiave}\\
\endfirsthead
\tabhead\hd{Rif.}&\hd{Parte}&\hd{Datasheet (sigla)}&\hd{Pagine chiave}\\
\endhead
\multicolumn{4}{r}{\footnotesize\itshape continua nella pagina seguente}\\
\endfoot
\bottomrule
\endlastfoot
U301 & STM32L562RET6P & ST DS12736 Rev 5 (DS-STM32L562) & p.\,33 TrustZone; p.\,65 VREFINT; p.\,87 pinout LQFP64 \enquote{P}; pp.\,94--126 tab.\,21 pin; p.\,126 note su PC13--15 e pin JTAG; pp.\,127--129 AF0--AF7; p.\,135 AF8--AF15; p.\,276 SPI\\
IC204 & Si4464-B1B-FMR & Silicon Labs Si4464/63/61/60 Rev 1.4 (DS-Si4464) & p.\,9 banda TX; p.\,13 abs.\ max; p.\,17 SPI; p.\,19 FRR; pp.\,20--21 stati, POR, SDN; p.\,35 PA\_PWR\_LVL; p.\,38 TCXO; pp.\,46--47 pin\\
Y201 & ECS-TXO-20CSMV4-320-AY & ECS (DS-ECS) & pp.\,1--2 codici di stabilità e temperatura\\
IC2 & QPA9510 & Qorvo, 23/06/2025 (DS-QPA9510); EVB 410--470\,MHz Test Summary 30/06/2025 (EVB Qorvo) & DS p.\,2 abs.\ max e tabella APC; p.\,3 caratteristiche; EVB pp.\,2--3 BOM, pp.\,6--7 guadagno e P1dB\\
IC202 & HMC349AMS8GE & ADI Rev.\ F (DS-HMC349A) & p.\,3 tempi; p.\,4 abs.\ max di potenza; p.\,8 tabella di verità e sequenza di accensione\\
Q201 & TQP3M9036 & Qorvo (DS-TQP3M9036) & p.\,2 shutdown, tempi, abs.\ max; p.\,7 dati a \SI{3.3}{\volt}\\
U701 & 74AUP2G14 & TI SN74AUP2G14 (DS-SN74AUP2G14) & p.\,1 $t_\text{pd}$\\
IC205 & TPS3825-33 & TI SLVS165 (DS-TPS3825) & pp.\,4--8 pin, soglie, $t_d$, MR\\
U403 & STWD100NYWY3F & ST DocID14134 Rev 8 locale, Rev 10 citata nel log (DS-STWD100) & Rev 8 p.\,14 trigger sui due fronti; Rev 10 tab.\,4 pp.\,12--13 (dal log)\\
U404 & DS3231M & Maxim (DS-DS3231M) & p.\,3 $t_\text{REC}$; p.\,7 pin; pp.\,11--15 registri; p.\,16 indirizzo\\
U401 & FM25V10 & Cypress 001-84499 Rev *D (DS-FM25V10) & pp.\,4--6 SPI e opcode; p.\,11 sleep e ID; p.\,17 $f_\text{SCK}$; p.\,19 $t_\text{PU}$\\
IC501 & TCAN3414 & TI (DS-TCAN3414) & p.\,1 velocità; p.\,3 pin; p.\,9 timeout di dominante\\
U402 & CH340E & WCH (DS-CH340E) & p.\,1 baud rate\\
IC4 & TPS2121 & TI SLVSEA3F (DS-TPS2121) & §7.5 VREF; §9 priorità\\
IC3 & TPSM863257 & TI (DS-TPSM863257) & p.\,7 soft-start\\
IC1, IC5 & TLV9051 & TI (DS-TLV9051) & offset (dal log R-18)\\
FL201 & LFCG-490+ & Mini-Circuits (DS-LFCG490) & perdita e potenza (dalla revisione del 04/11)\\
Q701 & FDN308P (clone TECH PUBLIC) & \textendash & solo funzione di interruttore (LOG P18)\\
\end{xltabular}}

\chapter{Da verificare e incongruenze}
\label{app:daverificare}

\section{Incongruenze tra schematico, log e datasheet}
Le incongruenze sono segnalate e non corrette: le decisioni spettano al team.

{\small
\rowcolors{2}{white}{sbrow}
\begin{xltabular}{\linewidth}{@{}L{8mm}L{58mm}Y@{}}
\tabhead\hd{\#}&\hd{Incongruenza}&\hd{Dettaglio e proposta}\\
\endhead
I-1 & PC13, PC14, PC15 collegati al dock nello schematico attuale & Netlist del 04/11/2026 15:22: PC13 = \net{GPIOSX\_2} (J708.13), PC14 = \net{GPIO\_PPS} (J708.17), PC15 = \net{GPIOSX\_0} (J708.14). Nella netlist della rev.\ 1.45 erano scollegati; il log (R-12, S-09, P15) dice che non arrivano all'MCU. Modifica non registrata nel log. Definire funzione e direzione.\\
I-2 & Ritardo del supervisore & F-05 e la nota sullo schematico UHF dicono \SI{200}{\milli\second}; il datasheet dà \SI{120}{}/\SI{200}{}/\SI{300}{\milli\second} (p.\,8). Il firmware deve usare \SI{300}{\milli\second}.\\
I-3 & Testo di F-14 & F-14 chiede \enquote{PA spento/APC a zero} prima di commutare PC2, ma l'APC è comandata da PC2. Si può solo togliere il pilotaggio (Si4464 fuori da TX). Confermare l'interpretazione.\\
I-4 & Tempo di accensione/spegnimento del QPA9510 & La revisione del 04/11 (P04) riporta \SI{1.7}{\micro\second}; nel testo estratto del datasheet (p.\,3) la riga \enquote{Turn ON/OFF Time} riporta 175 e la riga successiva 1.7 (probabilmente la corrente in A alla potenza massima). Il valore va letto sul PDF. Non cambia la regola dei \SI{10}{\micro\second}.\\
I-5 & F-09 e F-16 duplicate & Le due note contengono la stessa regola su EN32KHZ. Si possono unire.\\
I-6 & Nota \enquote{filtro passa banda} sul foglio UHF & Il filtro RX è un passa-basso. Decisione del log: ignorare la nota (P11).\\
I-7 & LED D203 nel modo B & La nota sullo schematico dice \enquote{D203 acceso = radio attiva}; nel modo B non è più vero (LOG §46).\\
I-8 & Versione del datasheet STWD100 & Il log cita la Rev 10 (tab.\,4, pp.\,12--13); il file locale è la Rev 8. I valori di $t_{WD}$ e $t_{PW}$ sono presi dal log.\\
I-9 & Etichetta \enquote{5v\_TX\_EN} su PC4 & La tabella della rev.\ 1.17 del log indica PC4 come \enquote{5v\_TX\_EN}; oggi PC4 è \net{RF\_MR\_N}. Testo storico superato.\\
\end{xltabular}}

\section{Dati da verificare}
{\small
\rowcolors{2}{white}{sbrow}
\begin{xltabular}{\linewidth}{@{}L{8mm}L{58mm}Y@{}}
\tabhead\hd{\#}&\hd{Dato}&\hd{Dove verificarlo}\\
\endhead
V-1 & Codici dei comandi e delle proprietà del Si4464 (POWER\_UP, PART\_INFO, GET\_INT\_STATUS, START\_TX/RX, CHANGE\_STATE, READ\_CMD\_BUFF, PA\_PWR\_LVL, GLOBAL\_XO\_TUNE) e valori di default & API Silicon Labs per la revisione B1 (AN625) e file generati da WDS. Valori presi dalla documentazione API della revisione C2A.\\
V-2 & Default di PA\_PWR\_LVL = \code{0x7F} (massima potenza) & Idem. Se confermato, F-01 va applicata prima di ogni TX.\\
V-3 & Tempo di POR del Si4464 dopo SDN basso & AN625 / datasheet; oggi si usa il polling di CTS con timeout.\\
V-4 & Valore di PA\_PWR\_LVL per \SI{\leq31}{\dBm} all'antenna & Misura M-01 sul prototipo.\\
V-5 & Livelli, polarità e uso di PPS, GPIOSX\_0, GPIOSX\_2 & Specifica del dock e decisione del team (I-1).\\
V-6 & Pull interni di PC13--PC15 & Dipende da V-5.\\
V-7 & Modo dei GPIO dopo il reset (analogico) & Reference manual RM0438, registro GPIOx\_MODER.\\
V-8 & VREF+ collegato internamente a VDDA nel package LQFP64 & Datasheet DS12736 (tabella dei pin) o RM0438.\\
V-9 & Soglia di power-fail $V_\text{PF}$ del DS3231M & DS-DS3231M p.\,3 (tabella non leggibile nel testo estratto).\\
V-10 & Frequenza I2C2 a \SI{400}{\kilo\hertz} con pull-up da \SI{10}{\kilo\ohm} & Misura del tempo di salita sul prototipo.\\
V-11 & Tolleranza iniziale del TCXO a \SI{25}{\celsius} & DS-ECS p.\,1.\\
V-12 & Tempo di accensione del QPA9510 (I-4) & DS-QPA9510 p.\,3 sul PDF.\\
V-13 & Campo di temperatura del TCXO (\SI{-30}{\celsius}) rispetto al requisito di missione & Decisione del team (LOG R-07, REV-1104 P07).\\
V-14 & Comando per la portante continua (CW) di collaudo & WDS / API Si4464.\\
\end{xltabular}}

\section{Misure di collaudo ancora da fare}
M-01 (potenza e valore F-01), M-02 (adattamento del PA), M-03 (3V3 in TX e scelta del modo A/B), M-04 (watchdog), M-05 (armoniche), M-06 (potenza all'MMCX e corrente del PA). Le misure M-01, M-03 e M-05 possono cambiare costanti del firmware (F-01, \code{MODO\_A/B}, riduzione di potenza).

\end{document}
